% ####################################################################
% Basic configuration of Beamera class and Jagiellonian theme
% ####################################################################
\RequirePackage[l2tabu, orthodox]{nag}



\ifx\PresentationStyle\notset
  \def\PresentationStyle{dark}
\fi



% Options: t - align frame text to the top.
\documentclass[10pt,t]{beamer}
\mode<presentation>
\usetheme[style=\PresentationStyle]{jagiellonian}




% ------------------------------------------------------------------------------------
% Procesing configuration files of Jagiellonian theme located
% in the directory "preambule"
% ------------------------------------------------------------------------------------
% Configuration for polish language
% Need description
\usepackage[polish]{babel}
% Need description
\usepackage[MeX]{polski}



% ------------------------------
% Better support of polish chars in technical parts of PDF
% ------------------------------
\hypersetup{pdfencoding=auto,psdextra}

% Package "textpos" give as enviroment "textblock" which is very usefull in
% arranging text on slides.

% This is standard configuration of "textpos"
\usepackage[overlay,absolute]{textpos}

% If you need to see bounds of "textblock's" comment line above and uncomment
% one below.

% Caution! When showboxes option is on significant ammunt of space is add
% to the top of textblock and as such, everyting put in them gone down.
% We need to check how to remove this bug.

% \usepackage[showboxes,overlay,absolute]{textpos}



% Setting scale length for package "textpos"
\setlength{\TPHorizModule}{10mm}
\setlength{\TPVertModule}{\TPHorizModule}


% ---------------------------------------
% Packages written for lectures "Geometria 3D dla twórców gier wideo"
% ---------------------------------------
% \usepackage{./Geometry3DPackages/Geometry3D}
% \usepackage{./Geometry3DPackages/Geometry3DIndices}
% \usepackage{./Geometry3DPackages/Geometry3DTikZStyle}
% \usepackage{./ProgramowanieSymulacjiFizykiPaczki/ProgramowanieSymulacjiFizykiTikZStyle}
% \usepackage{./Geometry3DPackages/mathcommands}


% ---------------------------------------
% Importing TikZ's libraries
% ---------------------------------------
\usetikzlibrary{arrows.meta}
\usetikzlibrary{positioning}





% % Configuration package "bm" that need for making bold symbols
\newcommand{\bmmax}{0}
\newcommand{\hmmax}{0}
\usepackage{bm}




% ---------------------------------------
% Packages for scientific texts
% ---------------------------------------
% \let\lll\undefined  % Sometimes you must use this line to allow
% "amsmath" package to works with packages with packages for polish
% languge imported
% /preambul/LanguageSettings/JagiellonianPolishLanguageSettings.tex.
% This comments (probably) removes polish letter Ł.
\usepackage{amsmath}  % Packages from American Mathematical Society (AMS)
\usepackage{amssymb}
\usepackage{amscd}
\usepackage{amsthm}
\usepackage{siunitx}  % Package for typsetting SI units.
\usepackage{upgreek}  % Better looking greek letters.
% Example of using upgreek: pi = \uppi


\usepackage{calrsfs}  % Zmienia czcionkę kaligraficzną w \mathcal
% na ładniejszą. Może w innych miejscach robi to samo, ale o tym nic
% nie wiem.







% !!!!!!!!!!!!!!!!!!!!!!!!!!!!!!
% !!!!!!!!!!!!!!!!!!!!!!!!!!!!!!
% EVIL STUFF
\if\JUlogotitle1
\edef\LogoJUPath{LogoJU_\JUlogoLang/LogoJU_\JUlogoShape_\JUlogoColor.pdf}
\titlegraphic{\hfill\includegraphics[scale=0.22]
{./JagiellonianPictures/\LogoJUPath}}
\fi

% ---------------------------------------
% Commands for handling colors
% ---------------------------------------


% Command for setting normal text color for some text in math modestyle
% Text color depend on used style of Jagiellonian

% Beamer version of command
\newcommand{\TextWithNormalTextColor}[1]{%
  {\color{jNormalTextFGColor}
    \setbeamercolor{math text}{fg=jNormalTextFGColor} {#1}}
}

% Article and similar classes version of command
% \newcommand{\TextWithNormalTextColor}[1]{%
%   {\color{jNormalTextsFGColor} {#1}}
% }



% Beamer version of command
\newcommand{\NormalTextInMathMode}[1]{%
  {\color{jNormalTextFGColor}
    \setbeamercolor{math text}{fg=jNormalTextFGColor} \text{#1}}
}


% Article and similar classes version of command
% \newcommand{\NormalTextInMathMode}[1]{%
%   {\color{jNormalTextsFGColor} \text{#1}}
% }




% Command that sets color of some mathematical text to the same color
% that has normal text in header (?)

% Beamer version of the command
\newcommand{\MathTextFrametitleFGColor}[1]{%
  {\color{jFrametitleFGColor}
    \setbeamercolor{math text}{fg=jFrametitleFGColor} #1}
}

% Article and similar classes version of the command
% \newcommand{\MathTextWhiteColor}[1]{{\color{jFrametitleFGColor} #1}}





% Command for setting color of alert text for some text in math modestyle

% Beamer version of the command
\newcommand{\MathTextAlertColor}[1]{%
  {\color{jOrange} \setbeamercolor{math text}{fg=jOrange} #1}
}

% Article and similar classes version of the command
% \newcommand{\MathTextAlertColor}[1]{{\color{jOrange} #1}}





% Command that allow you to sets chosen color as the color of some text into
% math mode. Due to some nuances in the way that Beamer handle colors
% it not work in all cases. We hope that in the future we will improve it.

% Beamer version of the command
\newcommand{\SetMathTextColor}[2]{%
  {\color{#1} \setbeamercolor{math text}{fg=#1} #2}
}


% Article and similar classes version of the command
% \newcommand{\SetMathTextColor}[2]{{\color{#1} #2}}










% ---------------------------------------
% Commands for few special slides
% ---------------------------------------
\newcommand{\EndingSlide}[1]{%
  \begin{frame}[standout]

    \begingroup

    \color{jFrametitleFGColor}

    #1

    \endgroup

  \end{frame}
}










% ---------------------------------------
% Commands for setting background pictures for some slides
% ---------------------------------------
\newcommand{\TitleBackgroundPicture}
{./JagiellonianPictures/Backgrounds/LajkonikDark.png}
\newcommand{\SectionBackgroundPicture}
{./JagiellonianPictures/Backgrounds/LajkonikLight.png}



\newcommand{\TitleSlideWithPicture}{%
  \begingroup

  \usebackgroundtemplate{%
    \includegraphics[height=\paperheight]{\TitleBackgroundPicture}}

  \maketitle

  \endgroup
}





\newcommand{\SectionSlideWithPicture}[1]{%
  \begingroup

  \usebackgroundtemplate{%
    \includegraphics[height=\paperheight]{\SectionBackgroundPicture}}

  \setbeamercolor{titlelike}{fg=normal text.fg}

  \section{#1}

  \endgroup
}










% ---------------------------------------
% Commands for lectures "Geometria 3D dla twórców gier wideo"
% Polish version
% ---------------------------------------
% Komendy teraz wykomentowane były potrzebne, gdy loga były na niebieskim
% tle, nie na białym. A są na białym bo tego chcieli w biurze projektu.
% \newcommand{\FundingLogoWhitePicturePL}
% {./PresentationPictures/CommonPictures/logotypFundusze_biale_bez_tla2.pdf}
\newcommand{\FundingLogoColorPicturePL}
{./PresentationPictures/CommonPictures/European_Funds_color_PL.pdf}
% \newcommand{\EULogoWhitePicturePL}
% {./PresentationPictures/CommonPictures/logotypUE_biale_bez_tla2.pdf}
\newcommand{\EUSocialFundLogoColorPicturePL}
{./PresentationPictures/CommonPictures/EU_Social_Fund_color_PL.pdf}
% \newcommand{\ZintegrUJLogoWhitePicturePL}
% {./PresentationPictures/CommonPictures/zintegruj-logo-white.pdf}
\newcommand{\ZintegrUJLogoColorPicturePL}
{./PresentationPictures/CommonPictures/ZintegrUJ_color.pdf}
\newcommand{\JULogoColorPicturePL}
{./JagiellonianPictures/LogoJU_PL/LogoJU_A_color.pdf}





\newcommand{\GeometryThreeDSpecialBeginningSlidePL}{%
  \begin{frame}[standout]

    \begin{textblock}{11}(1,0.7)

      \begin{flushleft}

        \mdseries

        \footnotesize

        \color{jFrametitleFGColor}

        Materiał powstał w ramach projektu współfinansowanego ze środków
        Unii Europejskiej w ramach Europejskiego Funduszu Społecznego
        POWR.03.05.00-00-Z309/17-00.

      \end{flushleft}

    \end{textblock}





    \begin{textblock}{10}(0,2.2)

      \tikz \fill[color=jBackgroundStyleLight] (0,0) rectangle (12.8,-1.5);

    \end{textblock}


    \begin{textblock}{3.2}(1,2.45)

      \includegraphics[scale=0.3]{\FundingLogoColorPicturePL}

    \end{textblock}


    \begin{textblock}{2.5}(3.7,2.5)

      \includegraphics[scale=0.2]{\JULogoColorPicturePL}

    \end{textblock}


    \begin{textblock}{2.5}(6,2.4)

      \includegraphics[scale=0.1]{\ZintegrUJLogoColorPicturePL}

    \end{textblock}


    \begin{textblock}{4.2}(8.4,2.6)

      \includegraphics[scale=0.3]{\EUSocialFundLogoColorPicturePL}

    \end{textblock}

  \end{frame}
}



\newcommand{\GeometryThreeDTwoSpecialBeginningSlidesPL}{%
  \begin{frame}[standout]

    \begin{textblock}{11}(1,0.7)

      \begin{flushleft}

        \mdseries

        \footnotesize

        \color{jFrametitleFGColor}

        Materiał powstał w ramach projektu współfinansowanego ze środków
        Unii Europejskiej w ramach Europejskiego Funduszu Społecznego
        POWR.03.05.00-00-Z309/17-00.

      \end{flushleft}

    \end{textblock}





    \begin{textblock}{10}(0,2.2)

      \tikz \fill[color=jBackgroundStyleLight] (0,0) rectangle (12.8,-1.5);

    \end{textblock}


    \begin{textblock}{3.2}(1,2.45)

      \includegraphics[scale=0.3]{\FundingLogoColorPicturePL}

    \end{textblock}


    \begin{textblock}{2.5}(3.7,2.5)

      \includegraphics[scale=0.2]{\JULogoColorPicturePL}

    \end{textblock}


    \begin{textblock}{2.5}(6,2.4)

      \includegraphics[scale=0.1]{\ZintegrUJLogoColorPicturePL}

    \end{textblock}


    \begin{textblock}{4.2}(8.4,2.6)

      \includegraphics[scale=0.3]{\EUSocialFundLogoColorPicturePL}

    \end{textblock}

  \end{frame}





  \TitleSlideWithPicture
}



\newcommand{\GeometryThreeDSpecialEndingSlidePL}{%
  \begin{frame}[standout]

    \begin{textblock}{11}(1,0.7)

      \begin{flushleft}

        \mdseries

        \footnotesize

        \color{jFrametitleFGColor}

        Materiał powstał w ramach projektu współfinansowanego ze środków
        Unii Europejskiej w~ramach Europejskiego Funduszu Społecznego
        POWR.03.05.00-00-Z309/17-00.

      \end{flushleft}

    \end{textblock}





    \begin{textblock}{10}(0,2.2)

      \tikz \fill[color=jBackgroundStyleLight] (0,0) rectangle (12.8,-1.5);

    \end{textblock}


    \begin{textblock}{3.2}(1,2.45)

      \includegraphics[scale=0.3]{\FundingLogoColorPicturePL}

    \end{textblock}


    \begin{textblock}{2.5}(3.7,2.5)

      \includegraphics[scale=0.2]{\JULogoColorPicturePL}

    \end{textblock}


    \begin{textblock}{2.5}(6,2.4)

      \includegraphics[scale=0.1]{\ZintegrUJLogoColorPicturePL}

    \end{textblock}


    \begin{textblock}{4.2}(8.4,2.6)

      \includegraphics[scale=0.3]{\EUSocialFundLogoColorPicturePL}

    \end{textblock}





    \begin{textblock}{11}(1,4)

      \begin{flushleft}

        \mdseries

        \footnotesize

        \RaggedRight

        \color{jFrametitleFGColor}

        Treść niniejszego wykładu jest udostępniona na~licencji
        Creative Commons (\textsc{cc}), z~uzna\-niem autorstwa
        (\textsc{by}) oraz udostępnianiem na tych samych warunkach
        (\textsc{sa}). Rysunki i~wy\-kresy zawarte w~wykładzie są
        autorstwa dr.~hab.~Pawła Węgrzyna et~al. i~są dostępne
        na tej samej licencji, o~ile nie wskazano inaczej.
        W~prezentacji wykorzystano temat Beamera Jagiellonian,
        oparty na~temacie Metropolis Matthiasa Vogelgesanga,
        dostępnym na licencji \LaTeX{} Project Public License~1.3c
        pod adresem: \colorhref{https://github.com/matze/mtheme}
        {https://github.com/matze/mtheme}.

        Projekt typograficzny: Iwona Grabska-Gradzińska \\
        Skład: Kamil Ziemian;
        Korekta: Wojciech Palacz \\
        Modele: Dariusz Frymus, Kamil Nowakowski \\
        Rysunki i~wykresy: Kamil Ziemian, Paweł Węgrzyn, Wojciech Palacz

      \end{flushleft}

    \end{textblock}

  \end{frame}
}



\newcommand{\GeometryThreeDTwoSpecialEndingSlidesPL}[1]{%
  \begin{frame}[standout]


    \begin{textblock}{11}(1,0.7)

      \begin{flushleft}

        \mdseries

        \footnotesize

        \color{jFrametitleFGColor}

        Materiał powstał w ramach projektu współfinansowanego ze środków
        Unii Europejskiej w~ramach Europejskiego Funduszu Społecznego
        POWR.03.05.00-00-Z309/17-00.

      \end{flushleft}

    \end{textblock}





    \begin{textblock}{10}(0,2.2)

      \tikz \fill[color=jBackgroundStyleLight] (0,0) rectangle (12.8,-1.5);

    \end{textblock}


    \begin{textblock}{3.2}(1,2.45)

      \includegraphics[scale=0.3]{\FundingLogoColorPicturePL}

    \end{textblock}


    \begin{textblock}{2.5}(3.7,2.5)

      \includegraphics[scale=0.2]{\JULogoColorPicturePL}

    \end{textblock}


    \begin{textblock}{2.5}(6,2.4)

      \includegraphics[scale=0.1]{\ZintegrUJLogoColorPicturePL}

    \end{textblock}


    \begin{textblock}{4.2}(8.4,2.6)

      \includegraphics[scale=0.3]{\EUSocialFundLogoColorPicturePL}

    \end{textblock}





    \begin{textblock}{11}(1,4)

      \begin{flushleft}

        \mdseries

        \footnotesize

        \RaggedRight

        \color{jFrametitleFGColor}

        Treść niniejszego wykładu jest udostępniona na~licencji
        Creative Commons (\textsc{cc}), z~uzna\-niem autorstwa
        (\textsc{by}) oraz udostępnianiem na tych samych warunkach
        (\textsc{sa}). Rysunki i~wy\-kresy zawarte w~wykładzie są
        autorstwa dr.~hab.~Pawła Węgrzyna et~al. i~są dostępne
        na tej samej licencji, o~ile nie wskazano inaczej.
        W~prezentacji wykorzystano temat Beamera Jagiellonian,
        oparty na~temacie Metropolis Matthiasa Vogelgesanga,
        dostępnym na licencji \LaTeX{} Project Public License~1.3c
        pod adresem: \colorhref{https://github.com/matze/mtheme}
        {https://github.com/matze/mtheme}.

        Projekt typograficzny: Iwona Grabska-Gradzińska \\
        Skład: Kamil Ziemian;
        Korekta: Wojciech Palacz \\
        Modele: Dariusz Frymus, Kamil Nowakowski \\
        Rysunki i~wykresy: Kamil Ziemian, Paweł Węgrzyn, Wojciech Palacz

      \end{flushleft}

    \end{textblock}

  \end{frame}





  \begin{frame}[standout]

    \begingroup

    \color{jFrametitleFGColor}

    #1

    \endgroup

  \end{frame}
}



\newcommand{\GeometryThreeDSpecialEndingSlideVideoPL}{%
  \begin{frame}[standout]

    \begin{textblock}{11}(1,0.7)

      \begin{flushleft}

        \mdseries

        \footnotesize

        \color{jFrametitleFGColor}

        Materiał powstał w ramach projektu współfinansowanego ze środków
        Unii Europejskiej w~ramach Europejskiego Funduszu Społecznego
        POWR.03.05.00-00-Z309/17-00.

      \end{flushleft}

    \end{textblock}





    \begin{textblock}{10}(0,2.2)

      \tikz \fill[color=jBackgroundStyleLight] (0,0) rectangle (12.8,-1.5);

    \end{textblock}


    \begin{textblock}{3.2}(1,2.45)

      \includegraphics[scale=0.3]{\FundingLogoColorPicturePL}

    \end{textblock}


    \begin{textblock}{2.5}(3.7,2.5)

      \includegraphics[scale=0.2]{\JULogoColorPicturePL}

    \end{textblock}


    \begin{textblock}{2.5}(6,2.4)

      \includegraphics[scale=0.1]{\ZintegrUJLogoColorPicturePL}

    \end{textblock}


    \begin{textblock}{4.2}(8.4,2.6)

      \includegraphics[scale=0.3]{\EUSocialFundLogoColorPicturePL}

    \end{textblock}





    \begin{textblock}{11}(1,4)

      \begin{flushleft}

        \mdseries

        \footnotesize

        \RaggedRight

        \color{jFrametitleFGColor}

        Treść niniejszego wykładu jest udostępniona na~licencji
        Creative Commons (\textsc{cc}), z~uzna\-niem autorstwa
        (\textsc{by}) oraz udostępnianiem na tych samych warunkach
        (\textsc{sa}). Rysunki i~wy\-kresy zawarte w~wykładzie są
        autorstwa dr.~hab.~Pawła Węgrzyna et~al. i~są dostępne
        na tej samej licencji, o~ile nie wskazano inaczej.
        W~prezentacji wykorzystano temat Beamera Jagiellonian,
        oparty na~temacie Metropolis Matthiasa Vogelgesanga,
        dostępnym na licencji \LaTeX{} Project Public License~1.3c
        pod adresem: \colorhref{https://github.com/matze/mtheme}
        {https://github.com/matze/mtheme}.

        Projekt typograficzny: Iwona Grabska-Gradzińska;
        Skład: Kamil Ziemian \\
        Korekta: Wojciech Palacz;
        Modele: Dariusz Frymus, Kamil Nowakowski \\
        Rysunki i~wykresy: Kamil Ziemian, Paweł Węgrzyn, Wojciech Palacz \\
        Montaż: Agencja Filmowa Film \& Television Production~-- Zbigniew
        Masklak

      \end{flushleft}

    \end{textblock}

  \end{frame}
}





\newcommand{\GeometryThreeDTwoSpecialEndingSlidesVideoPL}[1]{%
  \begin{frame}[standout]

    \begin{textblock}{11}(1,0.7)

      \begin{flushleft}

        \mdseries

        \footnotesize

        \color{jFrametitleFGColor}

        Materiał powstał w ramach projektu współfinansowanego ze środków
        Unii Europejskiej w~ramach Europejskiego Funduszu Społecznego
        POWR.03.05.00-00-Z309/17-00.

      \end{flushleft}

    \end{textblock}





    \begin{textblock}{10}(0,2.2)

      \tikz \fill[color=jBackgroundStyleLight] (0,0) rectangle (12.8,-1.5);

    \end{textblock}


    \begin{textblock}{3.2}(1,2.45)

      \includegraphics[scale=0.3]{\FundingLogoColorPicturePL}

    \end{textblock}


    \begin{textblock}{2.5}(3.7,2.5)

      \includegraphics[scale=0.2]{\JULogoColorPicturePL}

    \end{textblock}


    \begin{textblock}{2.5}(6,2.4)

      \includegraphics[scale=0.1]{\ZintegrUJLogoColorPicturePL}

    \end{textblock}


    \begin{textblock}{4.2}(8.4,2.6)

      \includegraphics[scale=0.3]{\EUSocialFundLogoColorPicturePL}

    \end{textblock}





    \begin{textblock}{11}(1,4)

      \begin{flushleft}

        \mdseries

        \footnotesize

        \RaggedRight

        \color{jFrametitleFGColor}

        Treść niniejszego wykładu jest udostępniona na~licencji
        Creative Commons (\textsc{cc}), z~uzna\-niem autorstwa
        (\textsc{by}) oraz udostępnianiem na tych samych warunkach
        (\textsc{sa}). Rysunki i~wy\-kresy zawarte w~wykładzie są
        autorstwa dr.~hab.~Pawła Węgrzyna et~al. i~są dostępne
        na tej samej licencji, o~ile nie wskazano inaczej.
        W~prezentacji wykorzystano temat Beamera Jagiellonian,
        oparty na~temacie Metropolis Matthiasa Vogelgesanga,
        dostępnym na licencji \LaTeX{} Project Public License~1.3c
        pod adresem: \colorhref{https://github.com/matze/mtheme}
        {https://github.com/matze/mtheme}.

        Projekt typograficzny: Iwona Grabska-Gradzińska;
        Skład: Kamil Ziemian \\
        Korekta: Wojciech Palacz;
        Modele: Dariusz Frymus, Kamil Nowakowski \\
        Rysunki i~wykresy: Kamil Ziemian, Paweł Węgrzyn, Wojciech Palacz \\
        Montaż: Agencja Filmowa Film \& Television Production~-- Zbigniew
        Masklak

      \end{flushleft}

    \end{textblock}

  \end{frame}





  \begin{frame}[standout]


    \begingroup

    \color{jFrametitleFGColor}

    #1

    \endgroup

  \end{frame}
}










% ---------------------------------------
% Commands for lectures "Geometria 3D dla twórców gier wideo"
% English version
% ---------------------------------------
% \newcommand{\FundingLogoWhitePictureEN}
% {./PresentationPictures/CommonPictures/logotypFundusze_biale_bez_tla2.pdf}
\newcommand{\FundingLogoColorPictureEN}
{./PresentationPictures/CommonPictures/European_Funds_color_EN.pdf}
% \newcommand{\EULogoWhitePictureEN}
% {./PresentationPictures/CommonPictures/logotypUE_biale_bez_tla2.pdf}
\newcommand{\EUSocialFundLogoColorPictureEN}
{./PresentationPictures/CommonPictures/EU_Social_Fund_color_EN.pdf}
% \newcommand{\ZintegrUJLogoWhitePictureEN}
% {./PresentationPictures/CommonPictures/zintegruj-logo-white.pdf}
\newcommand{\ZintegrUJLogoColorPictureEN}
{./PresentationPictures/CommonPictures/ZintegrUJ_color.pdf}
\newcommand{\JULogoColorPictureEN}
{./JagiellonianPictures/LogoJU_EN/LogoJU_A_color.pdf}



\newcommand{\GeometryThreeDSpecialBeginningSlideEN}{%
  \begin{frame}[standout]

    \begin{textblock}{11}(1,0.7)

      \begin{flushleft}

        \mdseries

        \footnotesize

        \color{jFrametitleFGColor}

        This content was created as part of a project co-financed by the
        European Union within the framework of the European Social Fund
        POWR.03.05.00-00-Z309/17-00.

      \end{flushleft}

    \end{textblock}





    \begin{textblock}{10}(0,2.2)

      \tikz \fill[color=jBackgroundStyleLight] (0,0) rectangle (12.8,-1.5);

    \end{textblock}


    \begin{textblock}{3.2}(0.7,2.45)

      \includegraphics[scale=0.3]{\FundingLogoColorPictureEN}

    \end{textblock}


    \begin{textblock}{2.5}(4.15,2.5)

      \includegraphics[scale=0.2]{\JULogoColorPictureEN}

    \end{textblock}


    \begin{textblock}{2.5}(6.35,2.4)

      \includegraphics[scale=0.1]{\ZintegrUJLogoColorPictureEN}

    \end{textblock}


    \begin{textblock}{4.2}(8.4,2.6)

      \includegraphics[scale=0.3]{\EUSocialFundLogoColorPictureEN}

    \end{textblock}

  \end{frame}
}



\newcommand{\GeometryThreeDTwoSpecialBeginningSlidesEN}{%
  \begin{frame}[standout]

    \begin{textblock}{11}(1,0.7)

      \begin{flushleft}

        \mdseries

        \footnotesize

        \color{jFrametitleFGColor}

        This content was created as part of a project co-financed by the
        European Union within the framework of the European Social Fund
        POWR.03.05.00-00-Z309/17-00.

      \end{flushleft}

    \end{textblock}





    \begin{textblock}{10}(0,2.2)

      \tikz \fill[color=jBackgroundStyleLight] (0,0) rectangle (12.8,-1.5);

    \end{textblock}


    \begin{textblock}{3.2}(0.7,2.45)

      \includegraphics[scale=0.3]{\FundingLogoColorPictureEN}

    \end{textblock}


    \begin{textblock}{2.5}(4.15,2.5)

      \includegraphics[scale=0.2]{\JULogoColorPictureEN}

    \end{textblock}


    \begin{textblock}{2.5}(6.35,2.4)

      \includegraphics[scale=0.1]{\ZintegrUJLogoColorPictureEN}

    \end{textblock}


    \begin{textblock}{4.2}(8.4,2.6)

      \includegraphics[scale=0.3]{\EUSocialFundLogoColorPictureEN}

    \end{textblock}

  \end{frame}





  \TitleSlideWithPicture
}



\newcommand{\GeometryThreeDSpecialEndingSlideEN}{%
  \begin{frame}[standout]

    \begin{textblock}{11}(1,0.7)

      \begin{flushleft}

        \mdseries

        \footnotesize

        \color{jFrametitleFGColor}

        This content was created as part of a project co-financed by the
        European Union within the framework of the European Social Fund
        POWR.03.05.00-00-Z309/17-00.

      \end{flushleft}

    \end{textblock}





    \begin{textblock}{10}(0,2.2)

      \tikz \fill[color=jBackgroundStyleLight] (0,0) rectangle (12.8,-1.5);

    \end{textblock}


    \begin{textblock}{3.2}(0.7,2.45)

      \includegraphics[scale=0.3]{\FundingLogoColorPictureEN}

    \end{textblock}


    \begin{textblock}{2.5}(4.15,2.5)

      \includegraphics[scale=0.2]{\JULogoColorPictureEN}

    \end{textblock}


    \begin{textblock}{2.5}(6.35,2.4)

      \includegraphics[scale=0.1]{\ZintegrUJLogoColorPictureEN}

    \end{textblock}


    \begin{textblock}{4.2}(8.4,2.6)

      \includegraphics[scale=0.3]{\EUSocialFundLogoColorPictureEN}

    \end{textblock}





    \begin{textblock}{11}(1,4)

      \begin{flushleft}

        \mdseries

        \footnotesize

        \RaggedRight

        \color{jFrametitleFGColor}

        The content of this lecture is made available under a~Creative
        Commons licence (\textsc{cc}), giving the author the credits
        (\textsc{by}) and putting an obligation to share on the same terms
        (\textsc{sa}). Figures and diagrams included in the lecture are
        authored by Paweł Węgrzyn et~al., and are available under the same
        license unless indicated otherwise.\\ The presentation uses the
        Beamer Jagiellonian theme based on Matthias Vogelgesang’s
        Metropolis theme, available under license \LaTeX{} Project
        Public License~1.3c at: \colorhref{https://github.com/matze/mtheme}
        {https://github.com/matze/mtheme}.

        Typographic design: Iwona Grabska-Gradzińska \\
        \LaTeX{} Typesetting: Kamil Ziemian \\
        Proofreading: Wojciech Palacz,
        Monika Stawicka \\
        3D Models: Dariusz Frymus, Kamil Nowakowski \\
        Figures and charts: Kamil Ziemian, Paweł Węgrzyn, Wojciech Palacz

      \end{flushleft}

    \end{textblock}

  \end{frame}
}



\newcommand{\GeometryThreeDTwoSpecialEndingSlidesEN}[1]{%
  \begin{frame}[standout]


    \begin{textblock}{11}(1,0.7)

      \begin{flushleft}

        \mdseries

        \footnotesize

        \color{jFrametitleFGColor}

        This content was created as part of a project co-financed by the
        European Union within the framework of the European Social Fund
        POWR.03.05.00-00-Z309/17-00.

      \end{flushleft}

    \end{textblock}





    \begin{textblock}{10}(0,2.2)

      \tikz \fill[color=jBackgroundStyleLight] (0,0) rectangle (12.8,-1.5);

    \end{textblock}


    \begin{textblock}{3.2}(0.7,2.45)

      \includegraphics[scale=0.3]{\FundingLogoColorPictureEN}

    \end{textblock}


    \begin{textblock}{2.5}(4.15,2.5)

      \includegraphics[scale=0.2]{\JULogoColorPictureEN}

    \end{textblock}


    \begin{textblock}{2.5}(6.35,2.4)

      \includegraphics[scale=0.1]{\ZintegrUJLogoColorPictureEN}

    \end{textblock}


    \begin{textblock}{4.2}(8.4,2.6)

      \includegraphics[scale=0.3]{\EUSocialFundLogoColorPictureEN}

    \end{textblock}





    \begin{textblock}{11}(1,4)

      \begin{flushleft}

        \mdseries

        \footnotesize

        \RaggedRight

        \color{jFrametitleFGColor}

        The content of this lecture is made available under a~Creative
        Commons licence (\textsc{cc}), giving the author the credits
        (\textsc{by}) and putting an obligation to share on the same terms
        (\textsc{sa}). Figures and diagrams included in the lecture are
        authored by Paweł Węgrzyn et~al., and are available under the same
        license unless indicated otherwise.\\ The presentation uses the
        Beamer Jagiellonian theme based on Matthias Vogelgesang’s
        Metropolis theme, available under license \LaTeX{} Project
        Public License~1.3c at: \colorhref{https://github.com/matze/mtheme}
        {https://github.com/matze/mtheme}.

        Typographic design: Iwona Grabska-Gradzińska \\
        \LaTeX{} Typesetting: Kamil Ziemian \\
        Proofreading: Wojciech Palacz,
        Monika Stawicka \\
        3D Models: Dariusz Frymus, Kamil Nowakowski \\
        Figures and charts: Kamil Ziemian, Paweł Węgrzyn, Wojciech Palacz

      \end{flushleft}

    \end{textblock}

  \end{frame}





  \begin{frame}[standout]

    \begingroup

    \color{jFrametitleFGColor}

    #1

    \endgroup

  \end{frame}
}



\newcommand{\GeometryThreeDSpecialEndingSlideVideoVerOneEN}{%
  \begin{frame}[standout]

    \begin{textblock}{11}(1,0.7)

      \begin{flushleft}

        \mdseries

        \footnotesize

        \color{jFrametitleFGColor}

        This content was created as part of a project co-financed by the
        European Union within the framework of the European Social Fund
        POWR.03.05.00-00-Z309/17-00.

      \end{flushleft}

    \end{textblock}





    \begin{textblock}{10}(0,2.2)

      \tikz \fill[color=jBackgroundStyleLight] (0,0) rectangle (12.8,-1.5);

    \end{textblock}


    \begin{textblock}{3.2}(0.7,2.45)

      \includegraphics[scale=0.3]{\FundingLogoColorPictureEN}

    \end{textblock}


    \begin{textblock}{2.5}(4.15,2.5)

      \includegraphics[scale=0.2]{\JULogoColorPictureEN}

    \end{textblock}


    \begin{textblock}{2.5}(6.35,2.4)

      \includegraphics[scale=0.1]{\ZintegrUJLogoColorPictureEN}

    \end{textblock}


    \begin{textblock}{4.2}(8.4,2.6)

      \includegraphics[scale=0.3]{\EUSocialFundLogoColorPictureEN}

    \end{textblock}





    \begin{textblock}{11}(1,4)

      \begin{flushleft}

        \mdseries

        \footnotesize

        \RaggedRight

        \color{jFrametitleFGColor}

        The content of this lecture is made available under a Creative
        Commons licence (\textsc{cc}), giving the author the credits
        (\textsc{by}) and putting an obligation to share on the same terms
        (\textsc{sa}). Figures and diagrams included in the lecture are
        authored by Paweł Węgrzyn et~al., and are available under the same
        license unless indicated otherwise.\\ The presentation uses the
        Beamer Jagiellonian theme based on Matthias Vogelgesang’s
        Metropolis theme, available under license \LaTeX{} Project
        Public License~1.3c at: \colorhref{https://github.com/matze/mtheme}
        {https://github.com/matze/mtheme}.

        Typographic design: Iwona Grabska-Gradzińska;
        \LaTeX{} Typesetting: Kamil Ziemian \\
        Proofreading: Wojciech Palacz,
        Monika Stawicka \\
        3D Models: Dariusz Frymus, Kamil Nowakowski \\
        Figures and charts: Kamil Ziemian, Paweł Węgrzyn, Wojciech
        Palacz \\
        Film editing: Agencja Filmowa Film \& Television Production~--
        Zbigniew Masklak

      \end{flushleft}

    \end{textblock}

  \end{frame}
}



\newcommand{\GeometryThreeDSpecialEndingSlideVideoVerTwoEN}{%
  \begin{frame}[standout]

    \begin{textblock}{11}(1,0.7)

      \begin{flushleft}

        \mdseries

        \footnotesize

        \color{jFrametitleFGColor}

        This content was created as part of a project co-financed by the
        European Union within the framework of the European Social Fund
        POWR.03.05.00-00-Z309/17-00.

      \end{flushleft}

    \end{textblock}





    \begin{textblock}{10}(0,2.2)

      \tikz \fill[color=jBackgroundStyleLight] (0,0) rectangle (12.8,-1.5);

    \end{textblock}


    \begin{textblock}{3.2}(0.7,2.45)

      \includegraphics[scale=0.3]{\FundingLogoColorPictureEN}

    \end{textblock}


    \begin{textblock}{2.5}(4.15,2.5)

      \includegraphics[scale=0.2]{\JULogoColorPictureEN}

    \end{textblock}


    \begin{textblock}{2.5}(6.35,2.4)

      \includegraphics[scale=0.1]{\ZintegrUJLogoColorPictureEN}

    \end{textblock}


    \begin{textblock}{4.2}(8.4,2.6)

      \includegraphics[scale=0.3]{\EUSocialFundLogoColorPictureEN}

    \end{textblock}





    \begin{textblock}{11}(1,4)

      \begin{flushleft}

        \mdseries

        \footnotesize

        \RaggedRight

        \color{jFrametitleFGColor}

        The content of this lecture is made available under a Creative
        Commons licence (\textsc{cc}), giving the author the credits
        (\textsc{by}) and putting an obligation to share on the same terms
        (\textsc{sa}). Figures and diagrams included in the lecture are
        authored by Paweł Węgrzyn et~al., and are available under the same
        license unless indicated otherwise.\\ The presentation uses the
        Beamer Jagiellonian theme based on Matthias Vogelgesang’s
        Metropolis theme, available under license \LaTeX{} Project
        Public License~1.3c at: \colorhref{https://github.com/matze/mtheme}
        {https://github.com/matze/mtheme}.

        Typographic design: Iwona Grabska-Gradzińska;
        \LaTeX{} Typesetting: Kamil Ziemian \\
        Proofreading: Wojciech Palacz,
        Monika Stawicka \\
        3D Models: Dariusz Frymus, Kamil Nowakowski \\
        Figures and charts: Kamil Ziemian, Paweł Węgrzyn, Wojciech
        Palacz \\
        Film editing: IMAVI -- Joanna Kozakiewicz, Krzysztof Magda, Nikodem
        Frodyma

      \end{flushleft}

    \end{textblock}

  \end{frame}
}



\newcommand{\GeometryThreeDSpecialEndingSlideVideoVerThreeEN}{%
  \begin{frame}[standout]

    \begin{textblock}{11}(1,0.7)

      \begin{flushleft}

        \mdseries

        \footnotesize

        \color{jFrametitleFGColor}

        This content was created as part of a project co-financed by the
        European Union within the framework of the European Social Fund
        POWR.03.05.00-00-Z309/17-00.

      \end{flushleft}

    \end{textblock}





    \begin{textblock}{10}(0,2.2)

      \tikz \fill[color=jBackgroundStyleLight] (0,0) rectangle (12.8,-1.5);

    \end{textblock}


    \begin{textblock}{3.2}(0.7,2.45)

      \includegraphics[scale=0.3]{\FundingLogoColorPictureEN}

    \end{textblock}


    \begin{textblock}{2.5}(4.15,2.5)

      \includegraphics[scale=0.2]{\JULogoColorPictureEN}

    \end{textblock}


    \begin{textblock}{2.5}(6.35,2.4)

      \includegraphics[scale=0.1]{\ZintegrUJLogoColorPictureEN}

    \end{textblock}


    \begin{textblock}{4.2}(8.4,2.6)

      \includegraphics[scale=0.3]{\EUSocialFundLogoColorPictureEN}

    \end{textblock}





    \begin{textblock}{11}(1,4)

      \begin{flushleft}

        \mdseries

        \footnotesize

        \RaggedRight

        \color{jFrametitleFGColor}

        The content of this lecture is made available under a Creative
        Commons licence (\textsc{cc}), giving the author the credits
        (\textsc{by}) and putting an obligation to share on the same terms
        (\textsc{sa}). Figures and diagrams included in the lecture are
        authored by Paweł Węgrzyn et~al., and are available under the same
        license unless indicated otherwise.\\ The presentation uses the
        Beamer Jagiellonian theme based on Matthias Vogelgesang’s
        Metropolis theme, available under license \LaTeX{} Project
        Public License~1.3c at: \colorhref{https://github.com/matze/mtheme}
        {https://github.com/matze/mtheme}.

        Typographic design: Iwona Grabska-Gradzińska;
        \LaTeX{} Typesetting: Kamil Ziemian \\
        Proofreading: Wojciech Palacz,
        Monika Stawicka \\
        3D Models: Dariusz Frymus, Kamil Nowakowski \\
        Figures and charts: Kamil Ziemian, Paweł Węgrzyn, Wojciech
        Palacz \\
        Film editing: Agencja Filmowa Film \& Television Production~--
        Zbigniew Masklak \\
        Film editing: IMAVI -- Joanna Kozakiewicz, Krzysztof Magda, Nikodem
        Frodyma

      \end{flushleft}

    \end{textblock}

  \end{frame}
}



\newcommand{\GeometryThreeDTwoSpecialEndingSlidesVideoVerOneEN}[1]{%
  \begin{frame}[standout]

    \begin{textblock}{11}(1,0.7)

      \begin{flushleft}

        \mdseries

        \footnotesize

        \color{jFrametitleFGColor}

        This content was created as part of a project co-financed by the
        European Union within the framework of the European Social Fund
        POWR.03.05.00-00-Z309/17-00.

      \end{flushleft}

    \end{textblock}





    \begin{textblock}{10}(0,2.2)

      \tikz \fill[color=jBackgroundStyleLight] (0,0) rectangle (12.8,-1.5);

    \end{textblock}


    \begin{textblock}{3.2}(0.7,2.45)

      \includegraphics[scale=0.3]{\FundingLogoColorPictureEN}

    \end{textblock}


    \begin{textblock}{2.5}(4.15,2.5)

      \includegraphics[scale=0.2]{\JULogoColorPictureEN}

    \end{textblock}


    \begin{textblock}{2.5}(6.35,2.4)

      \includegraphics[scale=0.1]{\ZintegrUJLogoColorPictureEN}

    \end{textblock}


    \begin{textblock}{4.2}(8.4,2.6)

      \includegraphics[scale=0.3]{\EUSocialFundLogoColorPictureEN}

    \end{textblock}





    \begin{textblock}{11}(1,4)

      \begin{flushleft}

        \mdseries

        \footnotesize

        \RaggedRight

        \color{jFrametitleFGColor}

        The content of this lecture is made available under a Creative
        Commons licence (\textsc{cc}), giving the author the credits
        (\textsc{by}) and putting an obligation to share on the same terms
        (\textsc{sa}). Figures and diagrams included in the lecture are
        authored by Paweł Węgrzyn et~al., and are available under the same
        license unless indicated otherwise.\\ The presentation uses the
        Beamer Jagiellonian theme based on Matthias Vogelgesang’s
        Metropolis theme, available under license \LaTeX{} Project
        Public License~1.3c at: \colorhref{https://github.com/matze/mtheme}
        {https://github.com/matze/mtheme}.

        Typographic design: Iwona Grabska-Gradzińska;
        \LaTeX{} Typesetting: Kamil Ziemian \\
        Proofreading: Wojciech Palacz,
        Monika Stawicka \\
        3D Models: Dariusz Frymus, Kamil Nowakowski \\
        Figures and charts: Kamil Ziemian, Paweł Węgrzyn,
        Wojciech Palacz \\
        Film editing: Agencja Filmowa Film \& Television Production~--
        Zbigniew Masklak

      \end{flushleft}

    \end{textblock}

  \end{frame}





  \begin{frame}[standout]


    \begingroup

    \color{jFrametitleFGColor}

    #1

    \endgroup

  \end{frame}
}



\newcommand{\GeometryThreeDTwoSpecialEndingSlidesVideoVerTwoEN}[1]{%
  \begin{frame}[standout]

    \begin{textblock}{11}(1,0.7)

      \begin{flushleft}

        \mdseries

        \footnotesize

        \color{jFrametitleFGColor}

        This content was created as part of a project co-financed by the
        European Union within the framework of the European Social Fund
        POWR.03.05.00-00-Z309/17-00.

      \end{flushleft}

    \end{textblock}





    \begin{textblock}{10}(0,2.2)

      \tikz \fill[color=jBackgroundStyleLight] (0,0) rectangle (12.8,-1.5);

    \end{textblock}


    \begin{textblock}{3.2}(0.7,2.45)

      \includegraphics[scale=0.3]{\FundingLogoColorPictureEN}

    \end{textblock}


    \begin{textblock}{2.5}(4.15,2.5)

      \includegraphics[scale=0.2]{\JULogoColorPictureEN}

    \end{textblock}


    \begin{textblock}{2.5}(6.35,2.4)

      \includegraphics[scale=0.1]{\ZintegrUJLogoColorPictureEN}

    \end{textblock}


    \begin{textblock}{4.2}(8.4,2.6)

      \includegraphics[scale=0.3]{\EUSocialFundLogoColorPictureEN}

    \end{textblock}





    \begin{textblock}{11}(1,4)

      \begin{flushleft}

        \mdseries

        \footnotesize

        \RaggedRight

        \color{jFrametitleFGColor}

        The content of this lecture is made available under a Creative
        Commons licence (\textsc{cc}), giving the author the credits
        (\textsc{by}) and putting an obligation to share on the same terms
        (\textsc{sa}). Figures and diagrams included in the lecture are
        authored by Paweł Węgrzyn et~al., and are available under the same
        license unless indicated otherwise.\\ The presentation uses the
        Beamer Jagiellonian theme based on Matthias Vogelgesang’s
        Metropolis theme, available under license \LaTeX{} Project
        Public License~1.3c at: \colorhref{https://github.com/matze/mtheme}
        {https://github.com/matze/mtheme}.

        Typographic design: Iwona Grabska-Gradzińska;
        \LaTeX{} Typesetting: Kamil Ziemian \\
        Proofreading: Wojciech Palacz,
        Monika Stawicka \\
        3D Models: Dariusz Frymus, Kamil Nowakowski \\
        Figures and charts: Kamil Ziemian, Paweł Węgrzyn,
        Wojciech Palacz \\
        Film editing: IMAVI -- Joanna Kozakiewicz, Krzysztof Magda, Nikodem
        Frodyma

      \end{flushleft}

    \end{textblock}

  \end{frame}





  \begin{frame}[standout]


    \begingroup

    \color{jFrametitleFGColor}

    #1

    \endgroup

  \end{frame}
}



\newcommand{\GeometryThreeDTwoSpecialEndingSlidesVideoVerThreeEN}[1]{%
  \begin{frame}[standout]

    \begin{textblock}{11}(1,0.7)

      \begin{flushleft}

        \mdseries

        \footnotesize

        \color{jFrametitleFGColor}

        This content was created as part of a project co-financed by the
        European Union within the framework of the European Social Fund
        POWR.03.05.00-00-Z309/17-00.

      \end{flushleft}

    \end{textblock}





    \begin{textblock}{10}(0,2.2)

      \tikz \fill[color=jBackgroundStyleLight] (0,0) rectangle (12.8,-1.5);

    \end{textblock}


    \begin{textblock}{3.2}(0.7,2.45)

      \includegraphics[scale=0.3]{\FundingLogoColorPictureEN}

    \end{textblock}


    \begin{textblock}{2.5}(4.15,2.5)

      \includegraphics[scale=0.2]{\JULogoColorPictureEN}

    \end{textblock}


    \begin{textblock}{2.5}(6.35,2.4)

      \includegraphics[scale=0.1]{\ZintegrUJLogoColorPictureEN}

    \end{textblock}


    \begin{textblock}{4.2}(8.4,2.6)

      \includegraphics[scale=0.3]{\EUSocialFundLogoColorPictureEN}

    \end{textblock}





    \begin{textblock}{11}(1,4)

      \begin{flushleft}

        \mdseries

        \footnotesize

        \RaggedRight

        \color{jFrametitleFGColor}

        The content of this lecture is made available under a Creative
        Commons licence (\textsc{cc}), giving the author the credits
        (\textsc{by}) and putting an obligation to share on the same terms
        (\textsc{sa}). Figures and diagrams included in the lecture are
        authored by Paweł Węgrzyn et~al., and are available under the same
        license unless indicated otherwise. \\ The presentation uses the
        Beamer Jagiellonian theme based on Matthias Vogelgesang’s
        Metropolis theme, available under license \LaTeX{} Project
        Public License~1.3c at: \colorhref{https://github.com/matze/mtheme}
        {https://github.com/matze/mtheme}.

        Typographic design: Iwona Grabska-Gradzińska;
        \LaTeX{} Typesetting: Kamil Ziemian \\
        Proofreading: Leszek Hadasz, Wojciech Palacz,
        Monika Stawicka \\
        3D Models: Dariusz Frymus, Kamil Nowakowski \\
        Figures and charts: Kamil Ziemian, Paweł Węgrzyn,
        Wojciech Palacz \\
        Film editing: Agencja Filmowa Film \& Television Production~--
        Zbigniew Masklak \\
        Film editing: IMAVI -- Joanna Kozakiewicz, Krzysztof Magda, Nikodem
        Frodyma


      \end{flushleft}

    \end{textblock}

  \end{frame}





  \begin{frame}[standout]


    \begingroup

    \color{jFrametitleFGColor}

    #1

    \endgroup

  \end{frame}
}











% ------------------------------------------------------
% BibLaTeX
% ------------------------------------------------------
% Package biblatex, with biber as its backend, allow us to handle
% bibliography entries that use Unicode symbols outside ASCII.
\usepackage[
language=polish,
backend=biber,
style=alphabetic,
url=false,
eprint=true,
]{biblatex}

\addbibresource{Systemy-operacyjne-Bibliography.bib}





% ------------------------------------------------------
% Importing packages, libraries and setting their configuration
% ------------------------------------------------------





% ------------------------------------------------------
% Local packages
% ------------------------------------------------------
% Local configuration of this particular presentation
\usepackage{./Local-packages/local-settings}










% ####################################################################
\title{Systemy operacyjne}
\subtitle{XX.~Program Awk}

\author{Kamil Ziemian \\
  \email}


% \date{}
% ####################################################################










% ####################################################################
% Beginning of the document
\begin{document}
% ####################################################################





% ################################################
% Number of chars: 50k+, 54k+,
% Text is adjusted to the left and words are broken at the end of the line.
\RaggedRight
% ################################################





% ################################################
\maketitle
% ################################################





% ######################################
\EndingSlide{Ta prezentacja wciąż nie jest ukończona.}
% ######################################





% ######################################
\begin{frame}
  \frametitle{Spis treści}


  \tableofcontents

\end{frame}
% ######################################










% ################################################
\section{Program Awk, wprowadzenie}
% ################################################


% ######################################
\begin{frame}
  \frametitle{Program Awk}


  Naukę \colorhref{https://en.wikipedia.org/wiki/AWK}{Awka} warto rozpocząć
  od wyjaśnienia jego osobliwej nazwy. Jest to akronim nazwisk jego twórców
  \colorhref{https://en.wikipedia.org/wiki/Alfred\_Aho}{Alfreda Aho},
  \colorhref{https://en.wikipedia.org/wiki/Peter\_J.\_Weinberger}
  {Petera Weinbergera}
  i~\colorhref{https://en.wikipedia.org/wiki/Brian\_Kernighan}
  {Briana Kernighana}. Pod koniec prezentacji, można znaleźć refleksje
  na temat Awka, autorstwa Kernighana, w~których przyznaje, że~wybór tej
  nazwy nie był najszczęśliwszy.

  By wyjaśnić do czego służy Awk, warto porównać go z~poleceniem
  \texttt{grep}. Jak już wiemy \texttt{grep} istnieje po to, byśmy mogli
  za jego pomocą przeszukiwać tekst, w~poszukiwaniu miejsc, które pasują do
  zadanego wzorca. Natomiast Awk służy do \alert{przetwarzania} tekstu, czy
  zmieniania go wedle podanych przez nas poleceń.

  Awka można również używać do poszukiwania fragmentów tekstu spełniających
  zadanych wzorzec, ale to nie jest jego główne zadanie. Awk potrafi
  przeszukiwać tekst w~poszukiwaniu zdanego wzorca, bo~jest to mu potrzebne
  do efektywnego przetwarzania tekstu, ale nie to jest sednem tego programu.

\end{frame}
% ######################################





% ######################################
\begin{frame}
  \frametitle{Korzystanie z~Awka}


  Pomimo tego, naukę Awka zaczniemy od wykorzystania go do przeszukiwania
  zawartości pewnego pliku tekstowego. Jest to uzasadnione tym, że~w~naszej
  ocenie jest to najprostsza droga do jego nauki.

  Zaczniemy od pracy z~plikiem \texttt{Items-to-sell-I.txt}. Zawiera on
  $51$ linii, zawierających informację o~rzeczach, jakie pewna firma ma do
  sprzedania. W~firmie tej panuje spory bałagan, toteż zawartość tego
  pliku też jest chaotyczna. Oto kilka jego pierwszych linii.

  \begin{tabular}{l l l l}
    Name & Type & Amount & Color \\
    Volvo & car & 10 & green \\
    Volkswagen & car & 20 & celeste \\
    Hisense & computers & 15 & silver \\
    CCM & bike & 20 & cyan \\
    Toyota & car & 7  & blue \\
    \hspace{2em} \vdots & & &
  \end{tabular}

\end{frame}
% ######################################





% ######################################
\begin{frame}
  \frametitle{Korzystanie z~Awka}


  \begin{tabular}{l l l l}
    Name & Type & Amount & Color \\
    Volvo & car & 10 & green \\
    Volkswagen & car & 20 & celeste \\
    Hisense & computers & 15 & silver \\
    \hspace{2em} \vdots & & &
  \end{tabular}

  Plik \texttt{Items-to-sell-I.txt} jest podzielony na cztery kolumny.
  Pierwsza linia pliku zawiera nagłówki poszczególnych kolumn, które
  opisują ich zawartość. Każda następna linia zawiera nazwę przedmiotu,
  jego typ, ilość sztuk do sprzedania i~kolor przedmiotu.

  Firma ta ma do sprzedania samochody, typ „car”, rower, typ „bike”,
  i~części do komputerów, typ „computers”. Jeśli dziwi Państwa taki
  zestaw przedmiotów do~sprzedania, to proszę pamiętać, że~w~firmie tej
  panuje duży bałagan. Dodajmy, że~„celeste” i~wszystkie inne nazwy
  kolorów, które pojawiają~się w~tym pliku, naprawdę istnieją w~języku
  angielskim.

\end{frame}
% ######################################





% ######################################
\begin{frame}
  \frametitle{Korzystanie z~Awka}


  Awka można używać do pracy z~dowolnym plikiem tekstowym, nie tylko
  takim, który składa~się z~linii podzielonych na ustaloną liczbę kolumn.
  Jednak prościej jest zacząć naukę od pracy z~plikiem, który jest równo
  podzielony na kolumny.

  By wypisać całą zawartość pliku na ekranie możemy użyć polecenia \\
  \texttt{\$ awk '\{ print \}' Items-to-sell-I.txt} \\
  \texttt{::>} \texttt{Name Type Amount Color} \\
  \texttt{::>} \texttt{Volvo car 10 green} \\
  \texttt{::>} \texttt{Volkswagen car 20 celeste} \\
  \texttt{::>} \texttt{Hisense computers 15 silver} \\
  \hspace{3em} \vdots

  Aby wypisać tylko pierwszą kolumnę możemy użyć polecenia: \\
  \texttt{awk '\{ print \$1 \}' Items-to-sell-I.txt}. \\
  Jak~się Państwo domyślają, symbol \$1 oznacza pierwszą kolumnę.

\end{frame}
% ######################################





% ######################################
\begin{frame}
  \frametitle{Korzystanie z~Awka}


  Wypisywanie pierwszej kolumny pliku. \\
  \texttt{\$ awk '\{ print \$1 \}' Items-to-sell-I.txt} \\
  \texttt{::>} \texttt{Name} \\
  \texttt{::>} \texttt{Volvo} \\
  \texttt{::>} \texttt{Volkswagen} \\
  \texttt{::>} \texttt{Hisense} \\
  \hspace{3em} \vdots

  Tak samo możemy wypisać tylko drugą kolumnę. \\
  \texttt{\$ awk '\{ print \$2 \}' Items-to-sell-I.txt} \\
  \texttt{::>} \texttt{Type} \\
  \texttt{::>} \texttt{car} \\
  \texttt{::>} \texttt{car} \\
  \texttt{::>} \texttt{computers} \\
  \hspace{3em} \vdots

\end{frame}
% ######################################





% ######################################
\begin{frame}
  \frametitle{Korzystanie z~Awka}


  Oto jak możemy wypisać kolumnę pierwszą i~drugą. \\
  \texttt{\$ awk '\{ print \$1, \$2 \}' Items-to-sell-I.txt} \\
  \texttt{::>} \texttt{Name Type} \\
  \texttt{::>} \texttt{Volvo car} \\
  \texttt{::>} \texttt{Volkswagen car} \\
  \texttt{::>} \texttt{Hisense computers} \\
  \hspace{3em} \vdots

  Jeśli w~ostatni poleceniu pominiemy przecinek pomiędzy symbolami kolumn
  „\texttt{\$1}” i~„\texttt{\$2}”, wynik będzie inny, niż byśmy pragnęli. \\
  \texttt{\$ awk '\{ print \$1 \$2 \}' Items-to-sell-I.txt} \\
  \texttt{::>} \texttt{NameType} \\
  \texttt{::>} \texttt{Volvocar} \\
  \texttt{::>} \texttt{Volkswagencar} \\
  \texttt{::>} \texttt{Hisensecomputers} \\
  \hspace{3em} \vdots

\end{frame}
% ######################################





% ######################################
\begin{frame}
  \frametitle{Korzystanie z~Awka}


  Czy \texttt{print} w~Awku można przekazać argument~\texttt{\$0}?
  Tak, argument ten oznacza całą linię tekstu. \\
  \texttt{\$ awk '\{ print \$0 \}' Items-to-sell-I.txt} \\
  \texttt{::>} \texttt{Name Type Amount Color} \\
  \texttt{::>} \texttt{Volvo car 10 green} \\
  \texttt{::>} \texttt{Volkswagen car 20 celeste} \\
  \texttt{::>} \texttt{Hisense computers 15 silver} \\
  \hspace{3em} \vdots

  Pomimo tego, że~w~ostatnim poleceniu równie dobrze mogło by być użyte
  samo \texttt{print}, to użycie argument \texttt{\$0} może niekiedy
  uprościć wiele rzeczy.

\end{frame}
% ######################################





% ######################################
\begin{frame}
  \frametitle{Filtrowanie danych}


  Oznaczenia kolumn są \alert{wyrażeniami}, co znaczy, iż są obliczane
  przez Awka. Dlatego, też \texttt{\$(1 + 1)} to dla niego to samo
  co~\texttt{\$2}, co widać w~poniższym przykładzie. \\
  \texttt{\$ awk '\{ print \$( 1 + 1 ) \}' Items-to-sell-I.txt} \\
  \texttt{::>} \texttt{Type} \\
  \texttt{::>} \texttt{car} \\
  \texttt{::>} \texttt{car} \\
  \texttt{::>} \texttt{computers} \\
  \hspace{3em} \vdots

  Przejdźmy teraz do ciekawszych możliwości Awka, mianowicie do~filtrowania
  danych na podstawie zadanego warunku. Aby zobaczyć tylko te linie, które
  zawierające informacje o~samochodach do sprzedania możemy użyć
  następującego polecenia. \\
  \texttt{awk '\$2 == "car" \hspace{0.25em} \{ print \}' Items-to-sell-I.txt}.

\end{frame}
% ######################################





% ######################################
\begin{frame}
  \frametitle{Filtrowanie danych}


  \texttt{\$ awk '\$2 == "car" \hspace{0.25em} \{ print \}'
    Items-to-sell-I.txt} \\
  \texttt{::>} \texttt{Volvo car 10 green} \\
  \texttt{::>} \texttt{Volkswagen car 20 celeste} \\
  \texttt{::>} \texttt{Toyota car 20 blue} \\
  \hspace{3em} \vdots

  Przyjrzymy~się następującej części tego polecenia:
  \texttt{'\$2 == "car" \hspace{0.25em} \{ print \}'}. Jego sens jest dość
  prosty, najpierw sprawdzamy, czy zawartość drugiej kolumny to słowo
  „car”. Jeśli jest to prawdą, to wykonujemy działanie \texttt{\{ print \}},
  czyli wypisujemy całą linię.

  Możemy teraz podać prosty przykład, wykorzystania Awka do przetwarzania
  tekst. Powiedzmy, że~chcemy utworzyć plik, który zawiera tylko informacje
  o~samochodach, które~są do sprzedania, gdyż informacje o~innych
  przedmiotach nie są nam obecnie potrzebne.

\end{frame}
% ######################################





% ######################################
\begin{frame}
  \frametitle{Filtrowanie danych}


  By~utworzyć taki plik z~informacjami o~samochodach, wystarczy, że~użyjemy
  następującego polecenie, korzystające z~przekierowania procesu
  na~plik. \\
  \texttt{\$ awk '\$2 == "car" \hspace{0.25em} \{ print \}'
    Items-to-sell-I.txt > } \\
  \texttt{Car-to-sell.txt}

  Nie powinno Państwa zaskoczyć, że~aby wypisać wszystkie linie, które
  \alert{nie} zawierają informacji o~samochodach do sprzedania, możemy
  użyć następującego polecenia. \\
  \texttt{\$ awk '\$2 != "car" \hspace{0.25em} \{ print \}'
    Items-to-sell-I.txt} \\
  \texttt{::>} \texttt{Name Type Amount Color} \\
  \texttt{::>} \texttt{Hisense computers 15 silver} \\
  \texttt{::>} \texttt{CCM bike 20 cyan} \\
  \texttt{::>} \texttt{Accell cabike 20 citron} \\
  \hspace{3em} \vdots

\end{frame}
% ######################################





% ######################################
\begin{frame}
  \frametitle{Filtrowanie danych}


  Widzimy, że~Awk zwrócił wypisał jedną linię, której nie powinien,
  mianowicie „Name Type Amount Color”. Awk zadziała poprawnie, bo
  rzeczywiście słowo „Type” jest różne od~„car”. To dość częsty przypadek,
  że~pierwsze linie plików zawierają nagłówek informujący o~jego
  zawartości, dlatego gdy zdobędziemy trochę więcej wiedzy, powiem co można
  zrobić, by tych linii nie dostawać.

  Rozpatrzmy inny przykład. Powiedzmy, że~z~jakiegoś powodu, chcemy
  zobaczyć informacje o~wszystkich przedmiotach, jakie ta firma ma na
  sprzedaż, które mają kolor niebieski. Jednocześnie, zupełnie nie
  interesuje nas, ile egzemplarzy każdego z~tych przedmiotów jest
  w~magazynie. By to zrobić, możemy użyć następującego polecenia.
  \texttt{\$ awk '\$4 == "blue" \hspace{0.25em} \{ print \$1, \$2, \$4 \}'
    Items-to-sell-I.txt} \\
  \texttt{::>} \texttt{Toyota car blue} \\
  \hspace{3em} \vdots

\end{frame}
% ######################################





% ######################################
\begin{frame}
  \frametitle{Filtrowanie danych}


  Dla lepszego zaznajomienia~się z~własnościami Awka, przerobimy
  jeszcze kilka przykładów. Oto jak wybrać wszystkie linie, opisujące
  przedmioty, których mamy $10$ egzemplarzy do sprzedania. \\
  \texttt{\$ awk '\$3 == 10  \{ print \}'
    Items-to-sell-I.txt} \\
  \texttt{::>} \texttt{Volvo car 10 green} \\
  \texttt{::>} \texttt{Abici bike 10 red} \\
  \hspace{3em} \vdots

  Ten przykład stanowi prostą ilustrację tego, że~wyrażenia „\$3 == 10”
  i~„10 == \$3” znaczą to samo. \\
  \texttt{\$ awk '10 == \$3 \{ print \}' Items-to-sell-I.txt} \\
  \texttt{::>} \texttt{Volvo car 10 green} \\
  \texttt{::>} \texttt{Abici bike 10 red} \\
  \hspace{3em} \vdots

\end{frame}
% ######################################





% ######################################
\begin{frame}
  \frametitle{Filtrowanie danych}


  Poszukajmy wszystkich przedmiotów, których liczba jest różna
  od~$10$. \\
  \texttt{\$ awk '\$3 != 10 \{ print \}' Items-to-sell-I.txt} \\
  \texttt{::>} \texttt{Name Type Amount Color} \\
  \texttt{::>} \texttt{Volkswagen car 20 celeste} \\
  \texttt{::>} \texttt{Hisense computers 15 silver} \\
  \hspace{3em} \vdots

  Ponownie Awk zwracam nam linię „Name Type Amount Color”, jeszcze chwilę
  nie będziemy~się tym przejmować.

  Oto jak wyszukujemy przedmioty, których liczba jest większe od $30$-tu. \\
  \texttt{\$ awk '\$3 > 30 \{ print \}' Items-to-sell-I.txt} \\
  \texttt{::>} \texttt{Name Type Amount Color} \\
  \texttt{::>} \texttt{Huawei computers 40 silver} \\
  \texttt{::>} \texttt{Gitane bike 45 dandelion} \\
  \hspace{3em} \vdots

\end{frame}
% ######################################





% ######################################
\begin{frame}
  \frametitle{Filtrowanie danych i~problemy z~nimi}


  Ten ostatni wynik jest dziwny. Czemu Awk stwierdza, że~prawdą jest
  „Amount > 30”? Przecież słowo „Amount” nie jest liczbą. Tak, ale
  w~komputerze wszystko jest ciągiem $0$ i~$1$, więc jeśli trzeba, wszystko
  może być zinterpretowane jako liczba. I~tutaj wchodzi Awk, który próbuje
  zinterpretować słowo „Amount” jako liczbę i~wychodzi mu, że jest ona
  większa od~$30$. Z~takimi dziwactwami trzeba~się nauczyć żyć, bo czemu
  życie miałoby być proste?

  Zróbmy jeszcze kilka przykładu. Wyszukajmy wszystkie przedmioty, których
  liczba jest większa, bądź równa $30$-tu. \\
  \texttt{\$ awk '\$3 >= 30 \{ print \}' Items-to-sell-I.txt} \\
  \texttt{::>} \texttt{Name Type Amount Color} \\
  \texttt{::>} \texttt{Huawei computers 40 silver} \\
  \texttt{::>} \texttt{Brennabor bike 30 red} \\
  \hspace{3em} \vdots

\end{frame}
% ######################################





% ######################################
\begin{frame}
  \frametitle{Filtrowanie danych}


  Poszukajmy teraz przedmiotów, których jest mniej niż~$10$. \\
  \texttt{\$ awk '\$3 < 10 \{ print \}' Items-to-sell-I.txt} \\
  \texttt{::>} \texttt{Toyota car 7 blue} \\
  \texttt{::>} \texttt{BH bike 9 blue} \\
  \hspace{3em} \vdots

  Poszukajmy teraz przedmiotów, których liczba jest mniejsza bądź
  równa~$10$. \\
  \texttt{\$ awk '\$3 <= 10 \{ print \}' Items-to-sell-I.txt} \\
  \texttt{::>} \texttt{Volvo car 10 green} \\
  \texttt{::>} \texttt{Toyota car 7 blue} \\
  \hspace{3em} \vdots

  W~obecnej chwili znamy już kilka podstawowych sposobów korzystania
  z~Awka. Zanim poznamy ich więcej, zatrzymamy~się na chwilę
  by~przyjrzeć~się kilku technicznym zagadnieniom Awka.

\end{frame}
% ######################################










% ################################################
\section{Trochę technicznych informacji o~Awku}
% ################################################


% ######################################
\begin{frame}
  \frametitle{Rekordy i~pola}


  Teraz będzie trochę teoretyczne nudzenie, jednak bez wprowadzenia
  pewnych pojęć, ciężko jest mówić o~Awku w~sposób zrozumiały dla jego
  doświadczonego użytkownika.

  Przypomnijmy jeszcze raz, że~program Awk służy do przetwarzanie tekstu.
  Jeden przykład tego już widzieliśmy, Awkowi można przekazać tekst,
  który składa~się z czterech kolumn, a~on stworzy tekst, składający~się
  tylko z~jego pierwszej kolumny.

  Tekst który Awk otrzymuje na wejście, jest dzielony przez niego
  na~\textbf{rekordy} (ang.~\textit{records}). Zwykle każdy rekord
  składa~się pojedynczej linii tekstu i~na tych zajęciach nie będziemy
  rozważali innych przypadków. Każdy rekord dzieli~się na pola
  (ang.~\textit{fields}). Przy domyślnych ustawieniach, każde osobne słowo
  jest uznane przez Awka za pole. O~tym jak zmienić sposób dzielenia
  rekordów na~pola, powiemy niedługo.

\end{frame}
% ######################################





% ######################################
\begin{frame}
  \frametitle{Działania i~wzorce w~Awku}


  Podstawowy sposób pracy Awka jest następujący. Tekst wejściowy dzielony
  jest na rekordy, następnie na każdym rekordzie wykonywane jest
  \textbf{działanie} (ang. \textit{action}), którego treść podajemy
  w~nawiasach wąsatych \texttt{\{ {} \}}.

  Wobec tego polecenie \\
  \texttt{awk '\{ print \}' Items-to-sell-I.txt} \\
  mówi, że~Awk ma podzielić tekst w~pliku \texttt{Items-to-sell-I.txt}
  na rekordy, następnie wykonać działanie \texttt{print} na każdym
  z~tych rekordów. W~tym wypadku sprowadza~się to do wypisania całego pliku
  na ekranie.

  Ogólna postać wyrażenia Awka jest następująca. \\
  \texttt{awk [opcje] 'wzorzec \{działanie\}' nazwa-pliku} \\
  Do tej pory nie używaliśmy Awka z~ustawionymi jawnie opcjami, za~to
  widzieliśmy już polecenie wykorzystujące odpowiedni \textbf{wzorzec}
  (ang.~\textit{patter}), jak choćby „\$3 == 10”.

\end{frame}
% ######################################





% ######################################
\begin{frame}
  \frametitle{Działania i~wzorce w~Awku}


  \texttt{awk '\$3 == 10 \{ print \}' Items-to-sell-I.txt} \\
  Z~polecenia powyżej korzystaliśmy już wcześniej, teraz zaś opiszemy jego
  działanie, korzystając z~poznanej wcześniej terminologii. Awk przechodzi
  tekst rekord po rekordzie (linia po linii), próbując dopasować do każdego
  rekordu wzorzec „\texttt{\$3 == 10}”. Jeśli ten wzorzec uda~się dopasować,
  czyli trzecia kolumna zawiera liczbę $10$, to Awk wykonuje na danym
  rekordzie działanie „\texttt{\{ print \}}”. Jeśli wzorzec nie zostanie
  dopasowany, to Awk nie robi nic i~przechodzi do następnego rekordu.

  Jak już~się Państwo zapewne domyślili, jeśli Awkowi nie prześlemy żadnego
  wzorca, jak w~poleceniu \\
  \texttt{awk '\{ print \}' Items-to-sell-I.txt} \\
  to uzna on, że~ma wykonać dane działanie na każdym rekordzie tekstu.

\end{frame}
% ######################################





% ######################################
\begin{frame}
  \frametitle{Ile tych Awków jest?}


  Jeśli masz system GNU/Linux, to na $99.999\%$ masz też do dyspozycji
  program Awk. Problem jest jednak taki, że~istnieje wiele nie do końca
  równoważnych wersji Awka. Dokładniej, Awk jest częścią standardu POSIX,
  można więc wziąć dokument tekstowy opisujący, co on ma robić i~samemu
  go~zaimplementować.

  Mówiąc inaczej, można wziąć ten dokument i~samemu napisać w~jakimś języku
  programowania, program realizujący zadane działania Awka. Wybór języka
  nie jest z~góry przewidziany, może to być C, C++, Go czy Python. Obecnie
  istnieją trzy różne implementacje Awka
  \colorhref{}{GNU Awk} (w~skrócie:
  \texttt{gawk}), \colorhref{}{\texttt{mawk}}
  i~\colorhref{}{\texttt{nawk}}. Niestety, nie
  każde polecenie, które działa dla \texttt{gawk}a, działa w~\texttt{mawk}u,
  etc. Jednak rdzeń zgodny z~\textsc{posix}em jest obecny w~nich wszystkich,
  co do pewnego stopnia łagodzi różnorakie wynikłe, z~braku pełnej
  kompatybilności tych trzech programów.


  Dwie podstawowe abstrakcje, których będziemy używać na tym kursie,
  to~\textbf{pliki} i~\textbf{procesy}. Nie wnikając w~szczegóły wystarczy
  nam stwierdzenie, że~plik czymś jest, a~proces coś robi. Warto tutaj
  przytoczyć następujące powiedzenie, które funkcjonuje w~wielu różnych
  wersjach i~różnych systemów operacyjnych.

  \textit{Wszystko w~systemie Linux jest plikiem. Jeśli coś nie jest
    plikiem to jest procesem.}

  Inaczej mówiąc, jeśli na komputerze znajdują~się jakieś dane, to są
  wewnątrz jakiegoś pliku. Jeśli komputer coś robi, to jest za to
  odpowiedzialny odpowiedni proces. Zrozumienie tego na tym poziomie jest
  wystarczające na~potrzeby tych zajęć. W~przyszłości może zajrzymy pod
  warstwę abstrakcji plików i~procesów, by~się przekonać, że~ich wewnętrzna
  budowa wcale nie jest banalna.

\end{frame}
% ######################################





% ######################################
\begin{frame}
  \frametitle{Ile tych Awków jest?}


  Polecenia lub programy \texttt{awk}a pozwalają nam przeszukiwać dane
  w~postaci tekstu, następnie zmieniać ten tekst na podstawie ustawionych
  przez nas parametrów i~własności.

  Dość szybko nasuwa~się nam pytanie: czym~się różni
  \colorhref{https://en.wikipedia.org/wiki/AWK}{\texttt{awk}}
  od~\colorhref{https://en.wikipedia.org/wiki/Grep}{\texttt{grep}a}?
  Zasadnicza sprowadza~się do tego, że~\texttt{grep} służy nam
  do~\alert{przeszukiwania} tekstu, podczas, gdy \texttt{awk} służy
  do \alert{przetwarzania} tekstu, a~przeszukiwanie tekstu jest tylko
  jedną z~jego funkcjonalności, która jest mu potrzebna by móc wykonywać
  swoje główne zadanie.

  Ze względu na to, że~\texttt{awk} jest pełnoprawnym językiem
  programowania, nie ma możliwości byśmy w~nauczyli~się wszystkich jego
  funkcjonalności w~czasie tego kursu. Będziemy celowali w~poznanie tych
  jego funkcjonalności, które są szczególnie przydatne w~pracy z~plikami
  tekstowymi. Nie należy jednak wynieść stąd błędnego wrażenia,
  że~\texttt{awk} nie jest dużym językiem programowania. Jest on
  stosunkowo mały, jednak wciąż na tyle duży, by móc pochłonąć trzy albo
  więcej spotkań.

  Warto wyjaśnić, skąd wzięła~się osobliwa nazwa tego programu. Jest to
  mianowicie akronim od nazwisk jego twórców:
  \colorhref{https://en.wikipedia.org/wiki/Alfred\_Aho}{Alfreda Aho}
  (ur.~$1941$),
  \colorhref{https://en.wikipedia.org/wiki/Peter\_J.\_Weinberger}
  {Petera Weinbergera} (ur.~$1942$)
  i~\colorhref{https://en.wikipedia.org/wiki/Brian\_Kernighan}
  {Briana Kernighana} (ur.~$1942$).

  \textit{I~have a natural revulsion to any operating system that shows so
    little plannning as to have to named all of its commands after digestive
    noises (awk, grep, fsck, nroff).} \\
    Autor nieznany, str.~$147$,
    \parencite{Garfinkel-Weise-Strassmann-The-UNIX-HATERS-Handbook-Pub-1994}.

\end{frame}
% ##################





% ##################
\begin{frame}
  \frametitle{Materiały do nauki \texttt{awk}a}



  Główny źródłem na którym oparta jest ta prezentacja tego polecenia/języka
  jest książka S.~Kenlona, D.~Morrissa i~R.~Younga
  \colorhref{https://opensource.com/downloads/awk-ebook}
  {\textit{A~practical guide to~learning \texttt{awk}}}
  \parencite{Kenlon-Morriss-Young-A-practical-guide-to-learning-awk-Ver-2025}.
  Niestety, ten e-book zawiera dużo błędów w~przykładach, pisałem w~tej
  sprawie do \href{opensource.com}{opensource.com}.

\end{frame}
% ######################################





% ######################################
\begin{frame}
  \frametitle{Korzystanie z~Awka}


  Przeczytaj pierwsze 5 stron. Strona 14.

  \texttt{awk '\{ print \$1 \}' Items-to-sell-I-CSV.txt} \\
  CSV comma separated value \\
  F from field separator \\
  \texttt{awk -F"," '\{ print \$1 \}' Items-to-sell-I-CSV.txt} \\
  \texttt{awk '{ print \$1 }' Items-to-sell-II.txt} \\
  \texttt{awk -F"," '\{ print \$1 \}' Items-to-sell-II.txt} \\
  \texttt{awk '\{ print \$1 \}' Items-to-sell-III.txt} \\
  \texttt{awk -F"," '\{ print \$1 \}' Items-to-sell-III.txt} \\
  \texttt{awk '\{ print NF ": " \$0 \}' Items-to-sell-II.txt} \\
  \texttt{awk -F '\{ print NF ": " \$0 \}' Items-to-sell-II.txt} \\
  \texttt{awk '/blue/ '\{ print \}' Items-to-sell-I.txt} \\
  Magia Awka \\
  \texttt{awk '\{ print > \$2 \}' Item-to-sell-I.txt} \\

\end{frame}
% ######################################





% ######################################
\begin{frame}
  \frametitle{Korzystanie z~Awka}


  \texttt{awk '\$1 $\sim$ /H.+a/ \{ print \$0 \}' Items-to-sell-I.txt} \\
  Zaczyna się na H i zawiera a. \\
  \texttt{patter \{ action \}} \\
  printf, \%d - decimal, \%s - string. \\
  \texttt{awk '\{ printf "\%d: \%s\textbackslash n", NR, \$0 \}'
    Items-to-sell-I.txt} \\
  \texttt{awk '\{ printf "Name: \%s; how many we have:
    \%d.\textbackslash n", \$1, \$3 \}' Items-to-sell-I.txt} \\
  \# - komentarz w Awku. \\
  \texttt{awk '\$3 > 30 \&\& \$2 == "car" \{ print \}' Items-to-sell-I.txt} \\
  Awk programs are essentially data-driven in that actions
  depend on the data, so they are quite a bit different from
  programs in many other programming languages. \\
  \texttt{awk '\$3 > 40 || \$2 == "bike" \{ print \}' Items-to-sell-I.txt} \\
  \texttt{awk 'NR > 1 \&\& (\$3 > 40 || \$2 == "bike") \{ print \}'
    Items-to-sell-I.txt} \\

  % Może już wysłałem te błędy.
  % Str. 9: printf -> {printf
  % Str. 11: You can use the logical operators and (written &&) and or
  % (written ||) to add specificity to your conditionals.
  % Bold symbols in wrong places.
  % Str. 15. Na samym początku, „egular” zamiast „regular”.

\end{frame}
% ######################################





% ######################################
\begin{frame}
  \frametitle{Czemu życie nie może być proste?}


  \texttt{awk -f ??add-stars-01.awk Items-to-sell-I.txt} \\
  \texttt{awk -f ???simple-use-of-variable-01.awk Items-to-sell-I.txt} \\
  \texttt{awk /Suzuki/ \{ print \} Items-to-sell-I.txt} \\
  \texttt{awk /uz/ \{ print \} Items-to-sell-I.txt} \\
  \texttt{awk /car/ \{ print \} Items-to-sell-I.txt} \\
  \texttt{awk /ca/ \{ print \} Items-to-sell-I.txt} \\
  \texttt{awk /ar/ \{ print \} Items-to-sell-I.txt} \\
  \texttt{awk \$4 $\sim$ /ar/ \{ print \} Items-to-sell-I.txt} \\
  \texttt{awk \$4 $\sim$ /blue/ \{ print \} Items-to-sell-I.txt} \\
  \texttt{awk \$4 !$\sim$ /blue/ \{ print \} Items-to-sell-I.txt} \\
  \texttt{awk /ol/ \{ print \} Items-to-sell-I.txt} \\
  \texttt{awk /o[lm]/ \{ print \} Items-to-sell-I.txt} \\
  \texttt{awk /d/ \{ print \} Items-to-sell-V.txt} \\
  \texttt{awk /\^{}d/ \{ print \} Items-to-sell-V.txt} \\

\end{frame}
% ######################################





% ##################
\begin{frame}
  \frametitle{????}


  \texttt{awk /d\$/ \{ print \} Items-to-sell-V.txt} \\
  \texttt{awk /t.b/ \{ print \} Items-to-sell-V.txt} \\
  % +, 1 albo więcej razy
  \texttt{awk /t.+b/ \{ print \} Items-to-sell-V.txt} \\
  % *, 0 albo więcej razy, czy to działa?
  \texttt{awk /t.*b/ \{ print \} Items-to-sell-V.txt} \\
  % ? to co przed nim, 0 lub 1 razy
  \texttt{awk /m?/ \{ print \} Items-to-sell-V.txt} \\
  % Substitute w całym rekordzie. Tylko pierwsze wystąpienie.
  \texttt{awk /m?/ \{ sub(/A1/, "D1"); print \} Items-to-sell-V.txt} \\
  % Substiture tylko w pierwszej kolumnie. Tylko pierwsze wystąpienie.
  \texttt{awk /m?/ \{ sub(/A1/, "D1", \$1); print \} Items-to-sell-V.txt} \\


  % Może już wysłałem te błędy.
  % Str. 9: printf -> {printf
  % Str. 11: You can use the logical operators and (written &&) and or
  % (written ||) to add specificity to your conditionals.
  % Bold symbols in wrong places.

\end{frame}
% ##################





% ######################################
\begin{frame}
  \frametitle{Działania i~wzorce w~Awku}


  Domyślnie, kązdy rekord (linia tekstu) dzielony jest na pola przez
  spacje i~inne białe znaki. Stąd każda linia w~pliku
  \texttt{Fruits-etc.txt} jest podzielona na trzy pola, które naturalnie
  ustawione są w~kolumny. Później powiem jak sprawić, by \texttt{awk} użył
  innych znaków jak separtorów pól (ang.~\textit{field separators}).

  Metodę pracy \texttt{awk}a możemy teraz opisać w~następujący sposób.
  Dzieli on dane wejściowe na linie (rekordy), następnie zaś na każdej
  z~linii wykonuje zadane działanie. Jak można już było zauważyć,
  \texttt{awk} przetwarza linie od góry do dołu. Przy czym składnia
  \texttt{awk}a jest taka, że~treść działania musi być ujęta w nawiasy
  wąsate \texttt{ \{ { } \} }.

\end{frame}
% ######################################




% ######################################
\begin{frame}
  \frametitle{Zmienna NR}


  Awk przechowuje numer rekordu w~zmiennej \texttt{NR}, więc jeśli chcemy
  zobaczyć wszystkie rekordy z~ich numerami, możemy skorzystać
  z~następującego polecenia. \\
  \texttt{\$ awk '\{ print NR ": " \$0 \}' Items-to-sell-I.txt} \\
  \texttt{::>} \texttt{1: Name Type Amount Color} \\
  \texttt{::>} \texttt{2: Volvo car 10 green} \\
  \texttt{::>} \texttt{3: Volkswagen car 20 celeste} \\
  \hspace{3em} \vdots

  Skorzystaliśmy tutaj też z~następującej własności \texttt{print} w~Awku.
  Mianowicie, jeśli prześlemy mu ciąg argumentów to on po prostu sklei je
  w~jeden tekst (string) i~go wyświetli. Dlatego działanie Awka
  „\texttt{print NR ": " \hspace{0.25em} \$0}” najpierw wypisuje numer
  rekordu, następnie dwukropek i~spację, a~potem zawartość rekordu (wiersza
  tekstu), którą reprezentuje zmienna~\texttt{\$0}.


\end{frame}
% ######################################





% ##################
\begin{frame}
  \frametitle{Jak rozsądnie korzystać z~\texttt{mv}?}


  Oznaczenia pól są \alert{wyrażeniami}, czyli \texttt{awk} oblicza
  wartość każdego z~ich wartość i~przekazuje ją dalej. Dletego też
  polecenie \\
  \texttt{\$ awk '\{ print \$(1 + 1) \}' Fruits-etc.txt} \\
  wypisze w~poprawny sposób drugą kolumnę pliku. Jak się komuś nudzi, to
  może spróbować następującego polecenia. \\
  \texttt{\$ awk '\{ print \$(sqrt(9)) \}' Fruits-etc.txt} \\
  Odradzamy używanie tak przekombinowanych rozwiązań, chyba że~dla własnej
  rozrywki.

  W~\texttt{awk}u możemy też kożystać z~dobrze  znanych w~informatyce
  operatorów porównania: \texttt{==}, \texttt{!=}, \texttt{>=}, \texttt{=<},
  \texttt{>}, \texttt{<}. W~większości wypadków, ich zrozumienie powinno być
  dość intuicyjne.

  Problem z~\colorhref{https://en.wikipedia.org/wiki/AWK}{\texttt{awk}iem}
  jest taki, że~jest to język programowania ukryty poleceniu wiersza
  poleceń. Mówiąc dokładniej \texttt{awk} jest językiem programowania
  specjalnego przeznaczenia. W~wypadu \texttt{awk}a dziedziną dla której
  został stworzony jest przetwarzanie plików tekstowych, niezależnie od
  tego, czy zawierają one wiersze Norwida (ja nie rozumiem wierszy
  Norwida), kod źródłowy w~języku~C, dane na temat cen surowców,~etc.

\end{frame}
% ##################









































% ######################################
\begin{frame}
  \frametitle{Korzystanie z~Awk}




\end{frame}
% ######################################








% ######################################
\section{Korzystanie z~\texttt{awk}a}
% ######################################














% ##################
\begin{frame}
  \frametitle{Filtrowanie danych}


  Jeśli przykładowo chcemy zobaczyć tylko te linię, które opisują
  owoc i~warzywa koloru żółtego, używamy polecenia \\
  \texttt{\$ awk '\$2 == "yellow" \{ print \$0 \}' Fruits-etc.txt} \\
  Jeśli chcemy zobaczyć wszystkie owoce i~warzywa, które \alert{nie} są
  żółte, to użyjemy polecenia \\
  \texttt{\$ awk '\$2 != "yellow" \{ print \$0 \}' Fruits-etc.txt} \\
  Ostatnie polecenie wypisze również, poprawnie rozumiejąc treść polecenia,
  również linię \\
  \begin{tabular}{l l l}
    name & color & amount \\
  \end{tabular} \\
  Oczywiście, można zaprogramować \texttt{awk}a, by nie wyświetlił nam
  pierwszej linii, sposoby by to zrobić poznamy później.

  Należy sobie jednak postawić pytanie, czy jeśli wszystkim czego pragniemy
  jest usunięcie pierwszej linii, to czy zaprogramowanie takie zachowania
  \texttt{awk}a nie pochłonie więcej energii, niż unięcie jej ręcznie?

\end{frame}
% ##################





% ##################
\begin{frame}
  \frametitle{Filtrowanie danych}


  Jak bowiem wiadomo, jedno z~mądrości informatyków głosi \textit{nigdy
    nie trać pięciu minut na zadanie, na którego automatyzajcę możesz
    poświęcić osiem godzin.}

  Przyjrzyjmy~się jeszcze poniższemu poleceniu, które wybiera
  wszystkie owoce i~warzywa, których mamy pięć i~więcej. \\
  \texttt{\$ awk '\$3 >= 5 \{ print \$0 \}' Fruits-etc.txt } \\
  Jeśli chcemy zapisać je wszystkie do osobnego pliku, to możemy skorzystać
  z~przekierowań, o~których mówiliśmy już wcześniej. \\
  \texttt{\$ awk '\$3 >= 5 \{ print \$0 \}' Fruits-etc.txt >
    Yellow-things-to-eat.txt} \\
  Linie te są teraz zapisane w~pliku
  \texttt{Yellow-things-to-eat.txt}.

  Wszystkie te własności możemy oczywiście łączyć, w~najróżniejszy sposób.
  Przykładowo, poniższe polecenie zapisuje do~pliku nazwę i~liczbę
  wszystkich żółtych owoców i~warzyw. \\
  \texttt{\$ awk '\$2 == "yellow" \{ print \$1, \$3 \} Fruits-etc.txt}

\end{frame}
% ##################





% ##################
\begin{frame}
  \frametitle{Filtrowanie danych}


  Jak bowiem wiadomo, jedno z~mądrości informatyków głosi \textit{nigdy
    nie trać pięciu minut na zadanie, na którego automatyzajcę możesz
    poświęcić osiem godzin.}

  Przyjrzyjmy~się jeszcze poniższemu poleceniu, które wybiera
  wszystkie owoce i~warzywa, których mamy pięć i~więcej. \\
  \texttt{\$ awk '\$3 >= 5 \{ print \$0 \}' Fruits-etc.txt } \\
  Jeśli chcemy zapisać je wszystkie do osobnego pliku, to możemy skorzystać
  z~przekierowań, o~których mówiliśmy już wcześniej. \\
  \texttt{\$ awk '\$3 >= 5 \{ print \$0 \}' Fruits-etc.txt >
    \textbackslash} \\
  \texttt{> Yellow-things-to-eat.txt} \\
  Linie te są teraz zapisane w~pliku
  \texttt{Yellow-things-to-eat.txt}.

  Wszystkie omawiane własności możemy łączyć ze sobą w~najróżniejszych
  konfiguracjach, również tych niezbyt oczywistych na pierwszy rzut oka.

\end{frame}
% ##################





% ##################
\begin{frame}
  \frametitle{Filtrowanie danych}


  Przykładowo, poniższe polecenie zapisuje do~pliku nazwę i~liczbę
  wszystkich żółtych owoców i~warzyw. \\
  \texttt{\$ awk '\$2 == "yellow" \{ print \$1, \$3 \} Fruits-etc.txt
    > \textbackslash} \\
  \texttt{> Yellow-things-to-eat.txt}.

  \texttt{\$ awk -F"," \{ print \$0 \} Fruits.csv} \\
  W~pewnych implementacja równoważnie można napisać \\
  \texttt{\$ awk --field-separator"," \{ print \$0 \} Fruits.csv}

\end{frame}
% ##################










% ######################################
\section{Przeszukiwanie według wzorca}
% ######################################



% ##################
\begin{frame}
  \frametitle{Przeszukiwanie według wzorca}


  Na chwilę zajmiemy~się trochę bardziej teoretycznymi sprawiami
  związanymi z~\texttt{awk}iem. Podstawowa forma jego ma następującą
  postać: \\
  \texttt{\$ awk [opcje] 'wzorzec \{działanie\}' nazwa-pliku} \\
  Jak sama nazwa wskazuje „\texttt{opcje}” można pominąć. Nazwa
  „\texttt{wzorzec}” oznacza wzorzec, którego \texttt{awk} będzie
  poszukiwał. Jeśli \texttt{awk} znajdzie \texttt{wzorzec} w~danej linii
  (rekordzie), to wykona na tej linii czynności zdefiniowane w~poleceniu
  \texttt{działąnie}. Również \texttt{wzorzec} może zostać pominięty, wtedy
  \texttt{awk} wypisze  każdą linię tekstu zawartą w~pliku który mu
  podaliśmy („\texttt{nazwa-pliku}”).

  % By przesadnie nie komplikować rozważań, będziemy przyjmować,
  % iż~\texttt{awk} ma wypisać każdą linię do której dopasował
  % „\texttt{wzorzec}”, chyba, że~jawnie zostanie stwierdzone inaczej.

  Warto dodać, że~działanie jak \texttt{awk} ma wykonać na linii pasującej
  do wzorca, musi być umieszczone wewnątrz nawiasów wąsatych \texttt{\{ \}}.
  Ich brak spowoduje błąd składni (ang. \textit{syntax error})
  \texttt{awk}a. Co tu mówić, jego informacje o~błędach nie są specjalnie
  pomocne, więc unikajmy ich jak tylko możemy.


  Separator pól .


  \texttt{\$ awk '/re/ \{ print \$0 \}' Fruits-etc.txt}
  Sterowane danymi \textit{data-driven}.

  Jak można~się z~tego domyślić, jeśli zdobędę dostęp do trybu jądra, mogę
  z~tym komputerem zrobić wszystko. Z~tego powodu hakerzy dokładają tylu
  starań, by uzyskać do niego dostęp \alert{bez} posiadania odpowiednich
  uprawnień. Dlatego też takie antycheaterskie programy takie jak Vanguard,
  który miał być dodany choćby do
  \colorhref{https://www.youtube.com/watch?v=nk6aKV2rY7E}{\textit{League~of
      Legends}}~są tak problematyczne
  \parencite{Low-Level-Why-Riots-anti-cheat-is-a-HUGE-problem-Ver-2024}.

  Vanguard działa bowiem w~trybie jądra, ma więc możliwość by zrobić
  z~moich komputerem absolutnie wszystko co chce. Czy ufamy firmie
  tworzącej ten program, że~nie nadużyje tych możliwości?

  Pytanie, czemu w~ogóle pozwalamy istnieć taki programom jak Vanguard?
  Dlatego, że~jeśli program ten pracuje w~trybie jądra to do wyszukiwania
  cheaterów używam zasobów \alert{naszych} komputera. Jeśli pracowałby
  w~trybie użytkownika, nie byłoby to możliwe, gdyż z~definicji trybu
  użytkownika, program ten nie miałby dostępu do informacji potrzebnych by
  wykryć cheatera.

\end{frame}
% ##################
















% % ######################################
% \section{Pliki i~procesy}
% % ######################################









% % ##################
% \begin{frame}
%   \frametitle{Pliki i~procesy}


%   Ponieważ pliki są mniej abstrakcyjne (hehe), zaczniemy nasze zajęcia od
%   zaznajomieniem~się z~pracą z~plikami pod powłoką \textsc{bash}.
%   Przyjmujemy przy tym, że~znają Państwo pewne pojęcia z~systemu Windows.

% \end{frame}
% % ##################










% % ######################################
% \section{BASH nie zna słowa „litość”. Część~I}
% % ######################################


% % ##################
% \begin{frame}
%   \frametitle{Litość? Co to jest?}


%   \textit{Gdy samochody rozwijały~się tak jak komputery, to dziś
%     Rolls-Royce kosztowałby jakieś 100\$, przejeżdżałaby milion mil na
%     jednym galonie benzyny i~raz do roku eksplodował zabijając wszystkich
%     w~środku.} \\
%   Robert Cringely, \textit{InfoWorld}, tłum. swobodne, str.~221,
%   \parencite{Garfinkel-Weise-Strassmann-The-UNIX-HATERS-Handbook-Pub-1994}.

%   Bóg jest litościwy, \textsc{bash} nie zna tego słowa. Przekonają~się
%   Państwo o~tym wielokrotnie. Każdy z~nas popełnia literówki i~\textsc{bash}
%   bez najmniejszej krztyny litości wytknie nam każdą z~nich, albo
%   z~jej powodu zinterpretuje wszystko na~opak. Jak z~wieloma innymi
%   rzeczami, trzeba nauczyć~się z~tym żyć.

%   Kolejny problem dotyczy białych znaków (ang. \textit{whitespaces}), przy
%   czym przez „białe znaki” rozumiemy wszystkie odstępy w~tekście, które na
%   kartce papieru są po prostu obszarami bieli (albo innego koloru papieru).
%   Trzy podstawowe białe znaki to spacja, tabulator i~znak nowej linii
%   (klawisz \texttt{Enter}).

% \end{frame}
% % ##################





% % ##################
% \begin{frame}
%   \frametitle{Swoboda używania białych znaków? Nie żartuj.}


%   Języki takie jak~C, czy nawet Python dają nam dużą swobodę
%   w~pisaniu białych znaków w~tekście. Wolność jaką mamy w~tym względzie
%   w~\textsc{bash}u, to wolność do podróżowania jaką mam więzień zamknięty
%   w~celi. Chodzi po celi można, ale czy to nazwiemy podróżą? Uczulam więc
%   Państwa, na to by pilnować, gdzie są białe znaki. Jeśli coś
%   w~\textsc{bash}u nie działa, to jest bardzo prawdopodobnie, iż~jest to
%   wynik spacji w~złym miejscu.

%   Będzie~się starali zaznaczać miejsca, gdzie szczególnie łatwo popełnić
%   błąd, jednak podanie wszystkich zasada pisania w~\textsc{bash}u i~wyjątków
%   od nich to naprawdę nie jest łatwa rzecz. Dlatego jeśli coś nie działa,
%   to proszę sprawdzić, czy wpisane polecenia wyglądają \alert{dokładnie}
%   tak jak zaprezentowane przykłady.

%   Błędy w~przykładach też mogą~się zdarzać, prosimy by pisać o~nich do nas
%   po adres \email.

% \end{frame}
% % ##################










% % ######################################
% \section{Podstawy pracy z~poziomu powłoki BASH}
% % ######################################



% % ##################
% \begin{frame}
%   \frametitle{Zaczynamy}


%   Zacznijmy od przyjęcia, że~pracujemy na komputerze jako użytkownik,
%   którego nazwa w~systemie operacyjnym została wybrana jako \texttt{adam}.
%   Jeśli włączymy konsolę zobaczymy pewnie coś takiego: \\
%   \texttt{adam:$\sim$\$} \\
%   Od tego momentu będziemy zwykle skracali ten znak zachęty \\
%   (ang.~\textit{prompt}) do symbolu \texttt{\$}. Przepisywanie pełnej jego
%   postaci wnosi bardzo mało informacji, za~to zabiera sporo miejsca.

%   Na początku zajmiemy~się tym, jak poruszać~się po systemie plików
%   i~katalogów. Na pewno wiedzą Państwo, że~na komputerze znajdują~się
%   katalogi (ang.~\textit{directories}), zwane też folderami, które zawierają
%   pliki i~inne katalogi. Z~punktu widzenia systemu GNU/Linux
%   katalog jest plikiem, jedynie specjalnego rodzaju. Dla tego systemu
%   wszystko jest plikiem lub procesem.

% \end{frame}
% % ##################





% % ##################
% \begin{frame}
%   \frametitle{Polecenia}


%   Podstawy działania katalogów w~systemie GNU/Linux są w~zasadzie takie
%   same jak w~systemie Windows, musimy~się jedynie przyzwyczaić do innych
%   konwencji używanych w~tym systemie. Dużo większym wyzwaniem zapewne
%   będzie zamienienie środowiska graficznego (\textsc{gui},
%   ang.~\textit{Graphic User Interface}) na pracę z~poziomu powłoki.

%   Podstawą pracy w~powłoce \textsc{bash} jest korzystanie z~poleceń (ang.
%   \textit{commands}). Bez wnikania w~szczegóły, każde polecenie posiada
%   nazwę i~wykonuje pewne działanie. By wykonać dane polecenie w~powłoce
%   należy najpierw wpisać jej nazwę, następnie zaś jej odpowiednie
%   opcje (flagi) i~argumenty. Część poleceń można wykonać bez wpisywania
%   opcji czy~argumentów.

%   Trzeba~się niestety przyzwyczaić do tego, że~nazwy poleceń, jak i~nazwy
%   ich opcji tych poleceń w~systemie GNU/Linux potrafią być dość dziwne
%   i~trudne do~zapamiętania.

% \end{frame}
% % ##################





% % ##################
% \begin{frame}
%   \frametitle{Polecenia}


%   Podstawowych poleceń systemu GNU/Linux jest około~$100$, do tego każde ma
%   kilkanaście albo więcej opcji. Z~tego powodu nauka ich jest bardzo
%   czasochłonna i~niełatwa. Teoretycznie w~każdej chwili w~\textsc{bash}u
%   możemy skorzystać z~manuala (podręcznika), gdzie są opisy wszystkich
%   poleceń. W~praktyce, żeby zrozumieć jego treść, wymagana jest już dość
%   spora wiedza.

%   Zdecydowaliśmy~się przyjąć następujące podejście. Naszym głównym celem
%   jest przekazanie Państwu umiejętności pracy spod powłoki \textsc{bash}
%   w~możliwie krótki i~bezbolesny sposób. Dlatego też, dość szybko
%   wprowadzimy możliwie dużo poleceń, bez zagłębiania~się w~zbyt wiele
%   szczegółów. Niekiedy będziemy wracać do wcześniej już omówionych poleceń,
%   by uzupełnić wiadomości o~nich. Wyczerpujące omówienie wszystkich
%   podstawowych poleceń i~tak zabrałoby zbyt dużo czasu, to zaś podejście
%   uważamy, za rozsądny kompromis.

% \end{frame}
% % ##################





% % ##################
% \begin{frame}
%   \frametitle{Polecenie \texttt{ls}}


%   Zacznijmy od~polecenia
%   \colorhref{https://linux.fandom.com/pl/wiki/Ls}{\texttt{ls}}. Jego nazwa
%   pochodzi prawdopodobnie od~ang.~\textit{List fileS}, ale są różne opinie
%   na ten temat. Polecenie to wyświetla zawartość katalogu w~którym~się
%   znajdujemy, z~tego powodu wynik działania tego polecenia będzie~się cały
%   czas zmieniać, jednym z~możliwych może być ten przedstawiony poniżej. \\
%   \texttt{\$ ls} \\
%   \texttt{{\color{jAxisBlue} Dokumenty}} \quad
%   \texttt{introduction-to-GNU-linux.pdf} \quad
%   \texttt{{\color{jAxisBlue} Pobrane}}

%   Polecenie \texttt{ls} jest przykładem takiego polecenia, które można
%   uruchomić bez dodatkowych opcji czy argumentów, ono samo ustali w~jakim
%   katalogu jesteśmy i~wypisze jego zawartość.

%   Obecnie jest standardem, że~powłoka używa kolorów i~różnych typów fontów,
%   by~informować nas o~typie i~własnościach pliku. W~przykładzie powyżej
%   nazwy katalogów są wyróżnione na niebiesko, trzeba jednak pamiętać,
%   że~wszystko to~zależy od~konkretnych ustawień systemu operacyjnego
%   i~łatwo może zostać zmienione.

% \end{frame}
% % ##################





% % ##################
% \begin{frame}
%   \frametitle{Polecenie \texttt{pwd}}


%   Kolejnym poleceniem jest
%   \colorhref{https://www.arturpyszczuk.pl/commands-pwd.html}{\texttt{pwd}}
%   od ang.~\textit{Print Working Directory}, rezultat jego działania można
%   zobaczyć poniżej. \\
%   \texttt{\$ pwd} \\
%   \texttt{/home/adam}

%   Ta krótka linia zawiera bardzo dużo informacji i~potrzebujemy przynajmniej
%   kilka z~nich wyjaśnić. Sposób w~jaki pliki i~katalogi są ułożone na~dysku
%   jest zdefiniowane przez \textbf{system plików} danego systemu
%   operacyjnego i~ten używany w~systemie GNU/Linux różni~się mocno tego
%   w~systemie Windows.

%   Zacznijmy od kilku uwag odnośnie notacji. Jeśli pracowalibyśmy w~tym
%   ostatnim systemie i~na dysku~C znajdowałby~się
%   katalog \texttt{Systemy-operacyjne-Ćwiczenia} a~w~nim katalog
%   \texttt{Zajęcia-01}, to ścieżka do niego miałaby postać: \\
%   \texttt{C:\textbackslash\textbackslash
%     Systemy-operacyjne-Ćwiczenia\textbackslash Zajęcia-01} \\

% \end{frame}
% % ##################





% % ##################
% \begin{frame}
%   \frametitle{System plików}


%   Odpowiednikiem ścieżki z~systemu Windows: \\
%   \texttt{C:\textbackslash\textbackslash
%     Systemy-operacyjne-Ćwiczenia\textbackslash Zajęcia-01} \\
%   w~systemie GNU/Linux byłaby ścieżka: \\
%   \texttt{/home/adam/Systemy-operacyjne-Ćwiczenia/Zajęcia-01} \\
%   Widzimy więc, że~gdy Windows do oddzielania nazw katalogów i~plików
%   używa symbolu „backslasha”: „\texttt{\textbackslash}”, to GNU/Linux
%   używa symbolu „slasha”: „\texttt{/}”. Jest to stosunkowo mała różnica,
%   ale trzeba~się do niej przyzwyczaić.

%   Poważniejsza różnica jest następująca. Komputery z~systemem GNU/Linux
%   mogą posiadać wiele dysków, nie jest to żadnym problemem, ale~system
%   plików GNU/Linuxa nie zna pojęcia dysku jako takiego. W~systemie GNU/Linux
%   istnienie naczelny katalog zwany rootem, czyli po polsku „korzeniem”,
%   w~którym zawierając~się wszystkie inne pliki oraz katalogi, nieważne na
%   którym fizycznym dysku są one umieszczone. Katalog root oznaczamy
%   symbolem~„\texttt{/}”.

% \end{frame}
% % ##################





% % % ##################
% % \begin{frame}
% %   \frametitle{System plików}


% %   Odpowiednikiem ścieżki z~systemu Windows: \\
% %   \texttt{C:\textbackslash\textbackslash
% %     Systemy-operacyjne-Ćwiczenia\textbackslash Zajęcia-01} \\
% %   w~systemie GNU/Linux byłaby ścieżka: \\
% %   \texttt{/home/adam/Systemy-operacyjne-Ćwiczenia/Zajęcia-01} \\
% %   Widzimy więc, że~gdy Windows do oddzielania nazw katalogów i~plików
% %   używa symbolu „backslasha”: „\texttt{\textbackslash}”, to GNU/Linux
% %   używa symbolu „slasha”: „\texttt{/}”. Jest to stosunkowo mała różnica,
% %   ale trzeba~się do niej przyzwyczaić.

% %   Poważniejsza różnica jest następująca. Komputery z~systemem GNU/Linux
% %   mogą posiadać wiele dysków, nie jest to żadnym problemem, ale~system
% %   plików GNU/Linuxa nie zna pojęcia dysku jako takiego. W~systemie GNU/Linux
% %   istnienie naczelny katalog zwany rootem, czyli po polsku „korzeniem”,
% %   w~którym zawierając~się wszystkie inne pliki oraz katalogi, nieważne na
% %   którym fizycznym dysku są one umieszczone. Katalog root oznaczamy
% %   symbolem~„\texttt{/}”.

% % \end{frame}
% % % ##################





% % ##################
% \begin{frame}
%   \frametitle{System plików}


%   Wróćmy do polecenia \\
%   \texttt{\$ pwd} \\
%   \texttt{/home/adam} \\
%   Linię \texttt{/home/adam} czytamy „W~katalogu root („\texttt{/}”)
%   znajduje~się katalog \texttt{home}, w~nim zaś znajduje~się katalog
%   \texttt{adam}.” Polecenie \texttt{pwd} podaje więc w~którym katalogu~się
%   obecnie znajdujemy, za~pomocą ścieżki, która zaczyna~się od katalogu root:
%   „\texttt{/}”.

%   Katalog w~którym obecnie jesteśmy nosi nazwę \textbf{katalogu roboczego}
%   (ang.~\textit{working directory}), a~jak powiedzieliśmy \texttt{pwd} jest
%   akronimem od~\textit{Print Working Directory}, czyli „wyświetl katalog
%   roboczy”.

%   Ścieżkę podającą lokalizację katalogu lub pliku nazywamy
%   \textbf{ścieżką bezwzględną}, jeśli zaczyna~się ona od katalogu root.
%   W~przeciwnym razie mówimy o~\textbf{ścieżce względnej} (podawanej
%   względem innego katalogu niż root).

% \end{frame}
% % ##################





% % ##################
% \begin{frame}
%   \frametitle{System plików}


%   Bez wielkich uproszczeń można przyjąć, że~\textsc{bash}u \alert{każdy}
%   ciąg symboli zaczynający od~symbolu „\texttt{/}” zostanie zinterpretowany
%   jako ścieżka do jakiegoś pliku lub katalogu. Z~tego powodu trzeba
%   \alert{bardzo} uważać przy stosowaniu symbolu~„\texttt{/}”.

%   W~szczególności, gdy stosujemy symbol „\texttt{/}”, trzeba zwracać dużą
%   uwagę na~spacje. Ścieżka postaci \\
%   \texttt{/home/adam} \\
%   jest poprawna. Ale już ścieżka \\
%   \texttt{/ home/adam} \\
%   jest \alert{błędna}! To w~jaki sposób jest błędna, to omówimy innym
%   razem. Doświadczenie uczy, że~na początku sprawia to ludziom
%   bardzo dużo problemów.

%   To, że~główny katalog nosi nazwę „korzenia” (ang.~\textit{root}) nie jest
%   przypadkiem, ale do tej kwestii wrócimy innym razem.

% \end{frame}
% % ##################





% % ##################
% \begin{frame}
%   \frametitle{Katalog domowy}


%   W~systemie GNU/Linux każdy użytkownik ma swój katalog domowy, w~którym
%   może przechowywać potrzebne mu pliki. Standardowo, choć mogą być wyjątki,
%   użytkownik \texttt{adam} otrzyma katalog \texttt{/home/adam}, użytkownik
%   \texttt{barbara} katalog \texttt{/home/barbara},~etc. Katalog domowy
%   bieżącego użytkownika posiada też alternatywną nazwę „$\sim$”. Stąd znak
%   zachęty \\
%   \texttt{adam:$\sim$\$} \\
%   przekazuje nam informację, że~działamy jako użytkownik \texttt{adam}
%   i~znajdujemy~się w~jego katalogu domowym.

%   Załóżmy, że~katalogu domowym znajduje~się katalog o~nazwie
%   \texttt{Systemy-operacyjne-Ćwiczenia}. By sprawdzić, czy tak jest
%   naprawdę możemy użyć polecenia \texttt{ls}: \\
%   \texttt{\$ ls} \\
%   \texttt{{\color{jAxisBlue} Dokumenty}} \quad
%   \texttt{introduction-to-GNU-linux.pdf} \quad
%   \texttt{{\color{jAxisBlue} Pobrane}} \\
%   \texttt{{\color{jAxisBlue} Systemy-operacyjne-Ćwiczenia}} \\

% \end{frame}
% % ##################





% % ##################
% \begin{frame}
%   \frametitle{Polecenie \texttt{cd}}


%   Jeśli w~bieżącym katalogu znajduje~się katalog
%   \texttt{Systemy-operacyjne-Ćwiczenia} to możemy do niego wejść
%   używając polecenia
%   \colorhref{https://linux.fandom.com/pl/wiki/Cd}{\texttt{cd}}
%   (od ang.~\textit{Change Directory}): \\
%   \texttt{adam:$\sim$\$ cd Systemy-operacyjne-Ćwiczenia} \\
%   \texttt{adam:$\sim$/Systemy-operacyjne-Ćwiczenia\$} \\
%   Wyjątkowo zapisaliśmy pełniejszą wersję znaku zachęty, by podkreślić,
%   iż~informuje on~nas w~którym katalogu jesteśmy. Mianowicie ścieżka
%   \texttt{$\sim$/Systemy-operacyjne-Ćwiczenia} informuje nas, że~w~katalogu
%   domowych użytkownika \texttt{adam} znajduje~się katalog
%   \texttt{Systemy-operacyjne-Ćwiczenia} i~w~nim właśnie jesteśmy.

%   Użyte wyżej polecenie \texttt{cd} jest przykładem polecenia, które
%   pobiera argument, mianowicie nazwę katalogu do którego możemy wejść.
%   W~istocie polecenie to pozwala nam na bardzo swobodne poruszanie~się
%   po systemie plików, jeśli opanujemy dobrze zapis ścieżek w~GNU/Linuxie,
%   ale wszystko po kolei.

% \end{frame}
% % ##################





% % ##################
% \begin{frame}
%   \frametitle{Polecenie \texttt{cd}}


%   Załóżmy, że~teraz chcielibyśmy przejść do katalogu nad nami. By to zrobić
%   należy przekazać poleceniu \texttt{cd} argument „\texttt{..}”: \\
%   \texttt{adam:$\sim$/Systemy-operacyjne-Ćwiczenia\$ cd ..} \\
%   \texttt{adam:$\sim$\$} \\
%   Dlaczego jako argument przesyłamy akurat „\texttt{..}” wyjaśnimy w~swoim
%   czasie, na razie przyjmijmy, iż~tak po prostu jest. Dla lepszej
%   ilustracji tego co~zrobiliśmy, ponownie wypisaliśmy znak zachęty,

%   Załóżmy, że~chcemy teraz wejść do katalogu
%   \texttt{Systemy-operacyjne-Ćwiczenia} i~utworzyć tam plik
%   \texttt{Notatki-z-zajęć.txt}. By utworzyć ten plik możemy skorzystać
%   z~polecenia
%   \colorhref{https://linux.fandom.com/pl/wiki/Touch}{\texttt{touch}}. Pełna
%   sekwencja poleceń będzie wyglądać tak. \\
%   \texttt{\$ cd Systemy-operacyjne-Ćwiczenia} \\
%   \texttt{\$ touch Notatki-z-zajęć.txt} \\

% \end{frame}
% % ##################





% % ##################
% \begin{frame}
%   \frametitle{Edycja plików z~pomocą programu gedit}


%   Możemy teraz użyć polecenia \texttt{ls} by~się przekonać, czy~plik
%   \texttt{Notatki-z-zajęć.txt} został poprawnie utworzony.
%   Z~komputerami nigdy nic nie wiadomo, więc warto to sprawdzić.

%   W~następnym kroku otworzymy go do edycji za~pomocą programu
%   \colorhref{https://gedit-text-editor.org/}{gedit}. By to zrobić
%   wpisujemy w~powłoce \\
%   \texttt{\$ gedit Notatki-z-zajęć.txt \&} \\
%   Proszę pamiętać o~dodaniu na końcu polecenia znaku
%   \alert{ampersand}:~„\&”. Inaczej gedit zablokuje nam możliwość wpisywania
%   kolejnych poleceń do powłoki i~by ją odzyskać będziemy musieli
%   wyłączyć gedita albo użyć kombinacji poleceń \texttt{jobs} i~\texttt{bg}.

%   Czemu uruchomienie gedita bez symbolu ampersanda ma takie skutki, omówimy
%   gdy będziemy dyskutowali zagadnienie procesów i~tego jak działają
%   w~systemie GNU/Linux. Nie jest to najprostszy temat na świecie.

% \end{frame}
% % ##################





% % ##################
% \begin{frame}
%   \frametitle{Okienko gedita}


%   Jeśli wszystko poszło dobrze, to na ekranie zobaczymy okienko Gedita,
%   które wygląda mniej więcej tak jak to na~zdjęciu~\eqref{fig:Gedit-window}.





%   \begin{figure}

%     \centering

%     \includegraphics[scale=0.145]
%     {./Presentations-pictures/Miscancellous-pictures/gedit-window.png}

%     \caption{Otwarte okno Gedita.}


%     \label{fig:Gedit-window}

%   \end{figure}

% \end{frame}
% % ##################





% % ##################
% \begin{frame}
%   \frametitle{Edytory tekstów}


%   Teraz możemy do pliku \texttt{Notatki-z-zajęć.txt} wpisać tekst, choćby
%   \textit{Notatki z~przedmiotu „Systemy operacyjne”}. Proszę pamiętać, by
%   zapisywać wprowadzone zmiany, służy do tego choćby skrót~\texttt{Ctrl-s}.

%   Kilka słów dlaczego wybraliśmy do~edycji plików program gedit. Na~tym
%   kursie korzystamy głównie z~niego, gdyż jest on prosty w~obsłudze
%   i~w~kilka sekund po jego otwarciu użytkownik może już pisać odpowiedni
%   tekst. Zostawiamy jednak Państwu pełną swobodę w~wyborze edytora tekstu.
%   Nie jest ważne jakiego programu~się używa, ważne jest, żeby udało~się
%   wykonać zadanie jakie przed nami stoi.

%   Jeśli chodzi o~edytory dostępne pod systemem GNU/Linux, to w~naszej
%   ocenie warto zwrócić uwagę choćby na~dwa owiane już legendą programy
%   do edycji tekstu \colorhref{https://www.vim.org/}{vim}
%   i~\colorhref{https://www.gnu.org/software/emacs/}{GNU Emacs}.

% \end{frame}
% % ##################










% % ######################################
% \section{Dalsze informacje o~pracy z~BASHem}
% % ######################################













% % ######################################
% \section{Autouzupełnianie w~BASHu}
% % ######################################


% % ##################
% \begin{frame}
%   \frametitle{Autouzupełnianie w~BASHu}


%   Autouzupełnianie to jedna z~funkcjonalności \textsc{bash}a, bez których
%   nie można w~nim normalnie pracować. Jest ono standardowo przypisane
%   klawiszowi \texttt{[TAB]}, jego działanie zilustrujemy na~przykładzie
%   polecenia~\texttt{mv}, ale w~pozostałych przypadkach działa ono
%   analogicznie.

%   Materiały pomocnicze do~tego zadania znajdują~się w~katalogu
%   \texttt{Materiały-do-prezentacji/SO-02-Podstawowe-ETC-Przykład/}.
%   W~katalogu tym znajdują~się pliki \texttt{Data-01.dat}
%   i~\texttt{Data-02.dat}, przy czym drugi z~nich chcemy przenieść
%   do~katalogu \texttt{Data-for-analysis/}.

%   Wprowadźmy w~powłoce \texttt{mv D} i~wciśnijmy klawisz \texttt{[TAB]}: \\
%   \texttt{\$ mv D [TAB]} \\
%   Na~podstawie wprowadzonej litery „\texttt{D}” i~znanej zawartości
%   katalogu powłoka postara~się uzupełnić jak największą część nazwy pliku,
%   dając w~rezultacie: \\
%   \texttt{\$ mv Data-}

% \end{frame}
% % ##################





% % ##################
% \begin{frame}
%   \frametitle{Autouzupełnianie w~BASHu}


%   Obecnie nasze polecenie ma postać \\
%   \texttt{\$ mv Data-} \\
%   W~tym momencie musimy wprowadzić dodatkowe znaki, tak by powłoka otrzymała
%   informację, czy chodzi o~plik \texttt{Data-01.dat},
%   \texttt{Data-02.dat}, czy katalog \texttt{Data-for-analysis/}. Po wpisaniu
%   znaku „\texttt{1}” możemy znów skorzystać z~autouzupełniania. \\
%   \texttt{\$ mv Data-02 [TAB]} \\
%   \texttt{\$ mv Data-02.dat}

%   Jeśli nie pamiętamy jak konkretnie nazywają~się pliki, które są w~danym
%   katalogu, to możemy kliknąć przycisk \texttt{[TAB]} dwa lub więcej razy,
%   wtedy powłoka pokaże nam nazwy wszystkich plików, które pasuję
%   do~dotychczas wprowadzonej nazwy. \\
%   \texttt{\$ mv Data- [TAB] \hspace{-1em} [TAB]} \\
%   \texttt{Data-01.dat} \quad \texttt{Data-02.dat} \quad \texttt{Data-for-analysis/}

% \end{frame}
% % ##################





% % ##################
% \begin{frame}
%   \frametitle{Autouzupełnianie w~BASHu}


%   Możemy zastosować autouzupełnianie również do nazwy katalogu do~którego
%   chcemy przenieść omawiany plik: \\
%   \texttt{\$ mv Data-02.dat Data-f [TAB]} \\
%   \texttt{\$ mv Data-02.dat Data-for-analysis/}

%   Praca w~powłoce bez autouzupełniana jest zupełnie nie do zniesienia,
%   dlatego polecamy Państwu potrenować użycie tej funkcjonalności w~praktyce.
%   Jak w~przypadku wielu innych rzecz, jest to najlepszy sposób by~się jej
%   nauczyć.

% \end{frame}
% % ##################










% % ######################################
% \section{BASH nie zna słowa „litość”. Część~II}
% % ######################################


% % ##################
% \begin{frame}
%   \frametitle{Litość? Jaka litość?}


%   Powłoka \textsc{bash} nie zna słowa litość. Powiedzmy, że~chcemy
%   zmienić nazwę pliku \texttt{Data-01.dat} na \texttt{Data-January-04.dat}
%   i~wpisujemy polecenie: \\
%   \texttt{\$ mv - i Data-01.dat Data-January-04.dat} \\
%   Wciskamy klawisz \texttt{Enter} i~widzimy coś takiego: \\
%   \texttt{mv: cel 'Data-January-04.dat nie jest katalogiem} \\
%   I~w~naszej głowie pojawia~się pytanie „Co~się dzieje?” Od razu zauważmy,
%   że~na Państwa komputerach treść otrzymanego komunikatu może być inna.
%   Problemowi okazuje~się być winna jest temu jedna, jedyna spacja między
%   symbolem „\texttt{-}” oraz~„\texttt{i}”.

%   Jeśli opcja zaczyna~się symbolem „\texttt{-}” to po nim
%   \alert{nie może być spacji}! Poprawna wersja tego polecenia to \\
%   \texttt{\$ mv -i Data-01.dat Data-January-04.dat}

%   Jak~się zapewne już Państwo domyślają, istnieją opcje poleceń nie
%   zaczynające~się od symbolu „\texttt{-}”, o~nich opowiemy sobie później.

% \end{frame}
% % ##################





% % ##################
% \begin{frame}
%   \frametitle{Logika BASHa? Jaka logika?}


%   Problem bierze~się stąd, że~jeśli wprowadzimy niepoprawne polecenie
%   \texttt{\$ mv - i Data-01.dat Data-January-04.dat} \\
%   to \textsc{bash} zinterpretuje „\texttt{-}”, „\texttt{i}” oraz
%   „\texttt{Data-01.dat}” jako nazwy \alert{plików}, które ma przenieść
%   do katalogu o~nazwie „\texttt{Data-January-04.dat}” i~zgłasza problem,
%   że~plik \texttt{Data-January-04.dat} nie jest katalogiem.

%   Polecenie \texttt{mv} zostało w~taki sposób napisane, że~oczekuje,
%   iż~przesłana do niej opcja zaczyna~się od symbolu „\texttt{-}”,
%   a~natychmiast po tym symbolu litera oznaczająca konkretną opcję.
%   W~przeciwnym razie zinterpretuje następujące po nim napisy jako nazwy
%   odpowiednich plików lub katalogów i~spróbuje wykonać na nich odpowiednią
%   operację.

%   Jak~się zapewne już Państwo domyślają, istnieją opcje poleceń nie
%   zaczynające~się od symbolu „\texttt{-}”, o~nich opowiemy sobie później.

%   Musimy podkreślić, że~w~tym co napisaliśmy jest wiele uproszczeń,
%   ale~nie chcieliśmy wchodzić teraz w~zbyt wiele szczegółów.

%   % Proszę pamiętać, że~\texttt{mv} jest programem, napisanym prawdopodobnie
%   % w~języku~C, który jakiś programista kiedyś stworzył i~ten twórca
%   % zdecydował, że~tak ma działać

%   % Jak~się zapewne już Państwo domyślają, istnieją opcje poleceń nie
%   % zaczynające~się od symbolu „\texttt{-}”, o~nich opowiemy sobie później.

% \end{frame}
% % ##################





% % ##################
% \begin{frame}
%   \frametitle{Polecenia też ktoś napisał}


%   % Problem bierze~się stąd, że~jeśli wprowadzimy niepoprawne polecenie
%   % \texttt{\$ mv - i Data-01.dat Data-January-04.dat} \\
%   % to \textsc{bash} zinterpretuje „\texttt{-}”, „\texttt{i}” oraz
%   % „\texttt{Data-01.dat}” jako nazwy \alert{plików}, które ma przenieść
%   % do katalogu o~nazwie „\texttt{Data-January-04.dat}” i~zgłasza problem,
%   % że~plik \texttt{Data-January-04.dat} nie jest katalogiem.

%   % Polecenie \texttt{mv} zostało w~taki sposób napisane, że~oczekuje,
%   % iż~przesłana do niej opcja zaczyna~się od symbolu „\texttt{-}”,
%   % a~natychmiast po tym symbolu litera oznaczająca konkretną opcję.
%   % W~przeciwnym razie zinterpretuje następujące po nim napisy jako nazwy
%   % odpowiednich plików lub katalogów i~spróbuje wykonać na nich odpowiednią
%   % operację.

%   % Musimy podkreślić, że~w~tym co napisaliśmy jest wiele uproszczeń,
%   % ale~nie chcieliśmy wchodzić teraz w~zbyt wiele szczegółów.

%   Proszę pamiętać, że~polecenie \texttt{mv} jest programem, napisanym
%   prawdopodobnie w~języku~C, który jakiś programista kiedyś stworzył i~ten
%   człowiek zdecydował, że~to polecenie ma działać tak, a~nie inaczej.
%   Często nie należy~się doszukiwać zbyt wiele logiki w~działaniu poleceń,
%   one są takie, bo~tak po prostu wyszło.

%   Gdybyśmy dziś tworzyli jakiś system operacyjny i~jego powłokę „od zera”,
%   to pewnie byśmy zaprojektowali je inaczej. Trzeba wszak uczyć~się na
%   własnych błędach ;).

% \end{frame}
% % ##################










% ######################################
\section{Pewne funkcjonalności \texttt{awk}a}
% ######################################


% % ##################
% \begin{frame}
%   \frametitle{Konwencja zapisu nazwy katalogów}


%   Jak już widzieliśmy, jeśli chcemy przejść do~katalogu \texttt{Systemy-} \\
%   \texttt{operacyjne-Ćwiczenia} to możemy użyć jednej z~dwóch
%   równoważnych nazw tego katalogu. \\
%   \texttt{\$ cd Systemy-operacyjne-Ćwiczenia/} \\
%   \texttt{\$ cd Systemy-operacyjne-Ćwiczenia}

%   W~dalszym ciągu tych notatek będziemy preferowali nazwy katalogów
%   kończą~się symbolem „\texttt{/}”, gdyż przekazują one dodatkową
%   informację, że~dany obiekt jest katalogiem.

% \end{frame}
% % ##################










% ######################################
\appendix
% ######################################





% ######################################
\EndingSlide{Dziękuję! Pytania?}
% ######################################










% % ######################################
% \section{Dlaczego BASH~się tak nazywa}
% % ######################################



% % ##################
% \begin{frame}
%   \frametitle{Dlaczego BASH~się tak nazywa?}


%   Historia zaczyna~się
%   od~\colorhref{https://en.wikipedia.org/wiki/Stephen\_R.\_Bourne}
%   {Steve Bourne’a}, który w~$1979$~roku napisał powłokę, która w~naturalny
%   sposób została nazwana
%   \colorhref{https://en.wikipedia.org/wiki/Bourne\_shell}{powłoką Bourne’a}
%   (ang.~\textit{Bourne shell}). Bazując na niej
%   \colorhref{https://en.wikipedia.org/wiki/Brian\_Fox\_(programmer)}
%   {Brian Jhan Fox} napisał w~$1989$~roku \textit{Bourne Again SHell},
%   czyli właśnie
%   \colorhref{https://en.wikipedia.org/wiki/Bash\_(Unix\_shell)}
%   {\textsc{bash}a}.

%   Sama nazwa \textit{Bourne again shell} to również żart językowy. Nazwę tę
%   można bowiem przetłumaczyć jako „ponownie powłoka Borne’a”, jak
%   i~„odrodzona na nowo powłoka”. W~tym drugim przypadku mamy do czynienia
%   z~grą słowną z~terminem
%   \colorhref{https://billygraham.org/answers/what-is-your-definition-of-a-born-again-christian}
%   {\textit{being born again}}, który zyskał
%   duże znaczenie teologiczne wśród protestantów z~nurtu
%   \colorhref{https://pl.wikipedia.org/wiki/Ewangelikalizm}{ewangelikanizmu}.

% \end{frame}
% % ##################





% % ##################
% \jagiellonianendslide{?????}
% % ##################










% ####################################################################
% ####################################################################
% Bibliography

\printbibliography





% ####################################################################
% End of the document

\end{document}
